\documentclass[aps,pre,preprint,groupedaddress]{revtex4}
%\usepackage[pdftex]{graphicx}
\usepackage{graphicx}
\usepackage{amsmath}
\usepackage{bm}
\usepackage{epsf,epsfig,graphics}
\usepackage{float}
\usepackage{makeidx}
\usepackage{latexsym}
\usepackage{amssymb}
\usepackage{mathrsfs}
\usepackage{amsbsy}
\usepackage{subfigure}
%\usepackage[table,x11names]{xcolor}
%\usepackage{kpfonts}
%\usepackage{xeCJK}
\usepackage{CJKutf8}
\usepackage{color}
\usepackage{morefloats}
\newcommand*{\vv}[1]{\vec{\mkern0mu#1}}
\newcommand{\figwsL}{16cm}
\newcommand{\figws}{7.5cm}
\newcommand{\figwss}{5.cm}
\newcommand{\figw}{10cm}
\newcommand{\figww}{15cm}
%\usepackage{mathaccents}
\DeclareMathAlphabet\mathbfcal{OMS}{cmsy}{b}{n}

\begin{document}

\title{Linear Response Theory in Noisy Network Dynamics for Nonequilibrium Steady States}
\author{Rong-Chih Chang$^1$\begin{CJK*}{UTF8}{bsmi}(張容誌)\end{CJK*}, Pik-Yin Lai$^{1,2}$%\footnote{pylai@phy.ncu.edu.tw}
\begin{CJK*}{UTF8}{bsmi}(黎璧賢)\end{CJK*}}
\affiliation{$^1$Dept. of Physics and Center for Complex Systems, National Central University, Chung-Li District, Taoyuan City 320, Taiwan, R.O.C.
}
\affiliation{$^2$Physics Division, National Center for Theoretical Sciences, Taipei 10617, Taiwan.}

\begin{abstract}
    Linear response theory predicts how a system responds to small perturbations from information contained in its unperturbed state. Here we develop a computable linear-response framework for noisy network dynamics near stable fixed points. The system consists of coupled stochastic variables with intrinsic node dynamics, weighted and possibly directed interactions, and Gaussian white noise. Linearization around the noise-free fixed point yields a Jacobian matrix and a steady-state covariance matrix that satisfy a Lyapunov equation. We derive explicit response formulas for observables under perturbations of node parameters, connection weights, and noise strengths. As a central example, we focus on the equal-time covariance response $\langle \delta x_i \delta x_j \rangle$ and express the corresponding susceptibilities in terms of correlation functions evaluated in the unperturbed steady state. The framework applies to both equilibrium steady states and nonequilibrium steady states, in which detailed balance is generally broken by asymmetric couplings or heterogeneous noise strengths. We further extend the theory to time-dependent perturbation protocols and obtain transient response formulas for means, forces, and covariance dynamics. Langevin simulations on undirected and directed networks confirm the predicted covariance responses within the linear-response regime.
\end{abstract}

\maketitle

\section{Introduction}
\quad In recent decades, advances in experimental techniques and data acquisition have made it possible to record high-dimensional time-series data in a wide range of biological, physical, and engineered systems. Many such interacting systems can be described as noisy network dynamics \cite{strogatz2001exploring, Mark2003, Mark2018}, in which heterogeneous components interact through weighted and possibly directed connections while being subject to stochastic fluctuations. These systems naturally generate time-series data from which covariance can be measured directly. From a physical perspective, many living systems are maintained through continuous exchanges of energy, matter, and information with their surroundings and therefore typically operate far from thermodynamic equilibrium \cite{PhysRevLett.71.2401, 10.1143/PTPS.130.29, sekimoto2010stochastic, PhysRevE.96.032123, PhysRevE.96.012601, Shiraishi2023}. Examples include metabolic processes and molecular motors. In such nonequilibrium systems, energy is continuously consumed and entropy is produced, while detailed balance can be broken by asymmetric interactions or by heterogeneous noise sources acting on different components.\par

One important inverse problem in complex networks is network reconstruction from time-series data. In many situations, the detailed connectivity is not directly accessible, whereas the node dynamics can be measured. Previous studies have shown that fluctuations can encode structural information and that network connectivity can be inferred from network dynamics \cite{PhysRevLett.104.058701, ching2013, ching2015, ching2017, Lai2017, ChengLai2022}. These results highlight the close relation between fluctuations and interactions in networked systems. Reconstruction alone, however, does not answer a closely related and equally important question: once the parameters of a noisy network are changed, how can one predict the resulting response of the system from properties of its original steady state? In many situations, the microscopic dynamics may be complicated, yet the system remains close to a stable fixed point and its fluctuation statistics can still be measured. This motivates a response theory formulated directly for noisy network dynamics near stable fixed points, where one can ask how perturbations of node parameters, couplings, or noise strengths modify steady-state observables.\par

In statistical mechanics, linear response theory (LRT) provides a general framework for predicting a system's response to small perturbations from correlation functions measured in the unperturbed state \cite{RKubo_1966, 10.1063/5.0003135}. In equilibrium systems, this connection is expressed by the fluctuation-dissipation theorem (FDT). Many realistic networked systems, however, do not operate near thermodynamic equilibrium, but instead settle into nonequilibrium steady states (NESS). Although LRT has been extended to NESS in several theoretical frameworks \cite{Seifert_2010, Mehl_2010, PhysRevLett.117.180601, PhysRevLett.134.157101}, these approaches do not directly provide explicit matrix-level susceptibilities for large networked systems. Therefore, there is a need for a practical framework that can be computed directly from the unperturbed system and then tested against perturbation experiments or simulations. This limitation becomes even more evident when one wishes to describe not only the difference between two steady states, but also the transient relaxation induced by a time-dependent perturbation protocol. Moreover, in coupled dynamical systems, local perturbations generally do not remain local. Perturbations of node parameters or interaction strengths can modify fluctuations and correlations throughout the system. In practical terms, local changes such as node removal, link failure, or targeted interventions can induce nontrivial changes in global observables \cite{NetworkDisruption}. A predictive framework for such responses is therefore useful not only for understanding noisy interacting systems, but also for assessing the effects of controlled perturbations and disruptions.

In this work, we address this problem for network dynamics fluctuating near stable fixed points. Each node is driven by Gaussian white noise and interacts with other nodes through weighted directed couplings, while also possessing its own intrinsic dynamics. Linearizing the dynamics around the noise-free stationary state leads to a Jacobian matrix $\mathbf{Q}$ and a steady-state equal-time covariance matrix $\mathbf{K}_0 \equiv \langle \delta \vec{x}\,\delta \vec{x}^{\top} \rangle$ satisfying a Lyapunov equation. Within this linearized Ornstein-Uhlenbeck description, the steady-state distribution is known exactly as a multivariate Gaussian. This exact solvability allows the first-order response functions to be derived explicitly from unperturbed steady-state quantities. For the original nonlinear noisy network, the same formulas should therefore be understood as the corresponding local Gaussian approximation near the stable fixed point. Within this scope, we develop a linear-response framework for both steady-state and time-dependent perturbations. Our central focus is the response of $\mathbf{K}_0$, for which we derive explicit and computable formulas under perturbations of intrinsic node parameters, connection weights, and noise strengths. The framework applies to both equilibrium steady states and nonequilibrium steady states of the linearized dynamics, where detailed balance is generally broken by asymmetric couplings or heterogeneous noise amplitudes. In this sense, the contribution of this work is not a general reformulation of NESS response theory for arbitrary nonlinear stochastic systems, but an operational response framework for linearized noisy networks near fixed points. The responses of means and forces further show that the same formalism also captures transient network relaxation.\par

The paper is organized as follows. In Sec. II, we introduce the noisy network model and its linearization near a stable fixed point. In Sec. III, we establish a linear response framework for both equilibrium and NESS in the linearized Langevin dynamics. In Sec. IV, we derive the steady-state linear response of $\mathbf{K}_0$ for our noisy network model. In Sec. V, we compare the theoretical predictions with Langevin simulations for both equilibrium and NESS cases. In Sec. VI, we extend the formalism to time-dependent perturbations and discuss the corresponding dynamical response.

\section{Noisy network dynamics and linearization near a stable fixed point}
\quad We consider a network of $N$ nodes, where the state of node i is described by a continuous variable $x_i(t)$. The intrinsic dynamics of each node is governed by a function $f_i(x_i;r_i)$, where $r_i$ denotes a node-specific parameter. Nodes $i$ and $j$ are connected by directed and weighted links described by the matrix ${\bf W}$. The noisy network dynamics are governed by
\begin{equation}
    \dot{x}_i = f_i(x_i;r_i) + \sum_{j\neq i} W_{ij} h(x_i, x_j) + \eta_i(t),\;i=1,2,...,N,
    \label{noisy_network}
\end{equation}
where $\eta_i(t)$ is zero-mean Gaussian white noise acting on node $i$. In many situations, the weighted matrix takes the $W_{ij}=g_{ij}A_{ij}$, with adjacency matrix elements $A_{ij}=0$ or 1, and connection weights $g_{ij}$. The coupling function $h(x_i,x_j)$ describes the interaction between nodes $i$ and $j$, the adjacency matrix or the connection weights are in general asymmetric for directed networks. It is convenient to combine the intrinsic and coupling terms into a deterministic force $\vec{F}(\vec{x};\vec{p})$, where $\vec{p}$ denotes a set of control parameters, including the node parameters $r_i$ and the network connection weights $W_{ij}$. Then the dynamics can be written as
\begin{equation}
    \dot{\vec{x}} = \vec{F}(\vec{x};\vec{p}) + \vec{\eta}(t),
    \label{deterministic_force}
\end{equation}
with
\begin{equation}
    F_j(\vec{x}; r_j, {\bf W}) = f_j(x_j;r_j) + \sum_{k} W_{jk} h(x_j,x_k).
\end{equation}
Here the parameters are assumed to be externally tunable and time-independent. Such a formulation is applicable to a wide range of noisy network systems. For example, in neuroscience, $x_i(t)$ may represent the membrane potential and $W_{ij}$ the strength and direction of synaptic connections; in ecological systems, $x_i(t)$ may represent the population of species $i$, coupled to other species through interaction terms.\par
We assume that the system fluctuates around a noise-free solution $\vec{X}$, where $\vec{F}(\vec{X}) = 0$ in the presence of noises. To analyze fluctuations near a stable fixed point, we introduce the deviation
\begin{equation}
    \Delta \vec{x} \equiv \vec{x} - \vec{X},
\end{equation}
and linearize the deterministic force around $\vec{X}$. The Jacobian matrix
\begin{equation}
    \mathbf{Q} = \nabla \vec{F}\big|_{\vec{X}}
\end{equation}
encodes both the network topology and the local dynamics, and is asymmetric in general. Up to linear order in $\Delta \vec{x}$, the stochastic dynamics becomes
\begin{equation}
    \Delta \dot{\vec{x}} = {\bf Q}\Delta \vec{x} + \vec{\eta}(t),
    \label{linearized_dynamics}
\end{equation}
with the Jacobian matrix which depends on the parameters
\begin{equation}
    Q_{ij} = W_{ij} \partial_2 h(X_i,X_j) + \left[ f_i'(X+i;r_i) + \sum_m W_{im} \partial_1 h(X_i,X_m) \right] \delta_{ij}.
\end{equation}
We assume that the fixed point $\vec{X}$ is linearly stable, such that all eigenvalues of $\mathbf{Q}$ have non-positive real parts. The noise is taken to be additive Gaussian white noise with covariance matrix
\begin{equation}
    \langle \vec{\eta}(t) \vec{\eta}^\intercal(t') \rangle = \boldsymbol{\sigma} \delta (t-t'),\;\langle \vec{\eta}(t) \rangle = 0,
\end{equation}
where, $\boldsymbol{\sigma}$ is chosen as a diagonal matrix. If the system reaches a steady state, the equal-time covariance matrix ${\bf K}_0 \equiv \langle \delta \vec{x} \delta \vec{x}^\intercal \rangle$ satisfies the Lyapunov equation
\begin{equation}
    \boldsymbol{\sigma} = -\mathbf{Q K}_0 - \mathbf{K}_0 \mathbf{Q}^\intercal.
    \label{Lyapunov_eq}
\end{equation}
This relation directly links the steady-state fluctuation statistics to the linearized drift matrix and noise matrix, which can be treated as the network version of FDT, and it is valid both in and out of equilibrium. With a given $\mathbf{Q}$ and $\boldsymbol{\sigma}$, $\mathbf{K}_0$ can be obtained by solving the Lyapunov equation. The probability current associated with the linearized dynamics is
\begin{equation}
    \vec{J}(\vec{x},t; \vec{p}) = \mathbf{Q} \delta \vec{x} \psi(\vec{x},t;\vec{p}) - \frac{1}{2}\boldsymbol{\sigma} \nabla \psi(\vec{x},t; \vec{p}).
    \label{probability_current}
\end{equation} The steady-state probability distribution of the linearized multidimensional system obeys the Gaussian distribution
\begin{equation}
    \psi_{ss}(\delta \vec{x}; \vec{p}) = \frac{1}{(2\pi)^{N/2}\sqrt{|{\bf K}_0|}} e^{-\frac{1}{2} \delta \vec{x}^\intercal {\bf K}^{-1}_0 \delta \vec{x}},
    \label{steady-state prob}
\end{equation}
which can be expressed as the effective potential $\Phi \equiv \frac{1}{2} \delta \vec{x}^\intercal {\bf K}^{-1}_0 \delta \vec{x}$. Then, the ensemble average of an observable $B(\vec{x}; \vec{p})$ is given by $\langle B \rangle = \int B(\vec{x}; \vec{p}) \psi_{ss}(\vec{x}) d\vec{x}$.
For a steady state, the current satisfies $\nabla \cdot \vec{J}_{ss} = 0$. Substituting the Gaussian distribution (\ref{steady-state prob}) into (\ref{probability_current}) gives
\begin{equation}
    \vec{J}_{ss} = \left( \mathbf{Q} + \frac{1}{2}\boldsymbol{\sigma} \mathbf{K}_0^{-1} \right) \delta \vec{x} \psi_{ss}.
\end{equation}
Thus, the existence of a stationary Gaussian distribution does not by itself imply equilibrium. The steady-state condition only requires the current to be divergence-free, whereas detailed balance requires it to vanish pointwise. The equilibrium condition is therefore
\begin{equation}
    \mathbf{Q} + \frac{1}{2} \boldsymbol{\sigma} \mathbf{K}_0^{-1} = 0.
\end{equation}
Using (\ref{Lyapunov_eq}), this condition is equivalent to
\begin{equation}
    \mathbf{Q}\boldsymbol{\sigma} = \boldsymbol{\sigma}\mathbf{Q}^\intercal,
\end{equation}
that is, $\mathbf{Q}\boldsymbol{\sigma}$ must be symmetric. This includes the familiar case of symmetric $\mathbf{Q}$ with uniform diagonal noise, but it is more general and does not require $\mathbf{Q}$ itself to be symmetric. For diagonal $\boldsymbol{\sigma}$, the condition becomes
\begin{equation}
    Q_{ij} \sigma_{jj} = Q_{ji} \sigma_{ii}.
\end{equation}
When $\mathbf{Q}\boldsymbol{\sigma} \neq \boldsymbol{\sigma}\mathbf{Q}^\intercal$, the couplings are asymmetric and/or the noise strengths are nonuniform across components. In that case, detailed balance is broken and the system reaches a NESS with persistent probability current, even though the distribution remains stationary. This distinction can also be detected from the time-lagged covariance $\mathbf{K}_{\tau} \equiv \langle \delta \vec{x}(t+\tau) \delta \vec{x}^\intercal(t) \rangle$: for equilibrium cases such as $\mathbf{W}=\mathbf{W}^\intercal$ with uniform diagonal noise, one expects $(\mathbf{K}_{\tau})_{ij}=(\mathbf{K}_{\tau})_{ji}$, whereas asymmetric couplings or heterogeneous noise strengths generally lead to $(\mathbf{K}_{\tau})_{ij}\neq(\mathbf{K}_{\tau})_{ji}$, reflecting broken time-reversal symmetry. It may also be viewed as a network generalization of Brownian-ratchet-like dynamics \cite{10.1143/PTPS.130.17}, where asymmetric couplings and heterogeneous fluctuations can generate directed nonequilibrium responses even in the absence of detailed balance.
% In the examples studied below, we mainly use the linear intrinsic dynamics $f_i = r_i(1-x_i)$ and the Logistic one $f_i = r_i x_i(1-x_i)$, together with the coupling function $h(x_i,x_j) = x_j-x_i$, for which the nodes tend to synchronize. 

\section{Steady-state linear response in noisy network dynamics}
% \subsection{}
We now develop a steady-state linear-response framework for noisy network dynamics under parameter perturbations. We focus on systems that relax to a stable fixed point and fluctuate around it under Gaussian white noise. For the linearized Ornstein-Uhlenbeck process introduced in Sec. II, the response identity derived below gives the exact first-order susceptibility. When it is applied back to the original nonlinear network, it should be interpreted as the corresponding fixed-point approximation, consistent with the local Gaussian steady-state form discussed in Sec. II and Appendix~A. When the parameters are perturbed as $\vec{p} \rightarrow \vec{p} + \delta \vec{p}$, the steady-state average of an observable ${\bf B}(\vec{x}; \vec{p})$ changes accordingly. Our goal is to express the response of steady-state observables in terms of quantities measured in the unperturbed steady state. We first consider how the noise-free fixed point depends on the parameters. In Eq. (\ref{deterministic_force}), the noise-free fixed point satisfies
\begin{equation}
    \vec{F}(\vec{X}; \vec{p}) = 0.
\end{equation}
When the parameters change as $\vec{p} \rightarrow \vec{p} + \delta \vec{p}$, the fixed point shifts from $\vec{X}$ to $\vec{X}+\delta \vec{X}$. Expanding
\begin{equation}
    \vec{F}(\vec{X}+\delta \vec{X}; \vec{p}+ \delta \vec{p}) = 0
\end{equation}
to linear order gives
\begin{equation}
    \nabla \vec{F}\big|_{\vec{X}; \vec{p}} \delta \vec{X} + \nabla_{\vec{p}}\vec{F}\big|_{\vec{X}; \vec{p}} \delta \vec {p} = 0.
\end{equation}
With $\mathbf{Q}\equiv\nabla\vec{F}\big|_{\vec{X}}$ and assuming that $\mathbf{Q}^{-1}$ exists, one obtains
\begin{equation}
    \frac{\partial X_i}{\partial p_\alpha} = -\left(\mathbf{Q}^{-1} \nabla_{\vec{p}} \vec{F}\big|_{\vec{X}; \vec{p}}\right)_{i\alpha}.
\end{equation}
This relation shows that the shift of the fixed point induced by a parameter perturbation is controlled by the linearized dynamics through $\mathbf{Q}$. For an observable $\mathbf{B}(\vec{x}; \vec{p})$, the steady-state linear response is defined by
\begin{equation}
    \langle \mathbf{B} \rangle_{ss, \vec{p}+\delta \vec{p}} = \langle \mathbf{B} \rangle_{ss, \vec{p}} + \boldsymbol{\chi} \delta \vec{p}.
    \label{response}
\end{equation}
The steady-state response function is then
\begin{equation}
    \boldsymbol{\chi} = \langle \nabla_{\vec{p}} \mathbf{B} \rangle_{ss, \vec{p}} + \int d\vec{x} \mathbf{B} \nabla_{\vec{p}} \psi_{ss}.
    \label{response_function}
\end{equation}
The steady state distribution can be written as $\psi_{ss}=Z^{-1}e^{-\Phi}$, and we define the conjugate force with respect to the parameters by
\begin{equation}
    \vec{\mathscr{F}} \equiv -\nabla_{\vec{p}} \Phi.
    \label{conjugate_force}
\end{equation} 
Notice that $\Phi=\Phi(\vec{p}, \vec{X}(\vec{p}))$, where the fixed point depends on the parameters. One obtains
\begin{equation}
    \partial_{p_\alpha} \psi_{ss} = \psi_{ss}(\mathscr{F}_\alpha-\langle \mathscr{F}_\alpha \rangle).
\end{equation}
Therefore,
\begin{equation}
    \boldsymbol{\chi} = \langle \nabla_{\vec{p}} \mathbf{B} \rangle_{ss, \vec{p}} + \langle \delta \mathbf{B} \delta \vec{\mathscr{F}}^\intercal \rangle_{ss, \vec{p}}.
    \label{response_function2}
\end{equation}
This is the steady-state linear-response formula used throughout this work \cite{softmatter,Seifert_2010,Mehl_2010}. It requires a differentiable steady-state distribution near a stable fixed point, but it does not require equilibrium. In the present paper, the explicit formulas derived from it are evaluated within the linearized Gaussian description, where they give exact first-order susceptibilities for the Ornstein-Uhlenbeck dynamics and leading-order approximations for the original nonlinear system near the fixed point. In the following sections, we apply this framework to perturbations of the intrinsic node parameter $r_\alpha$, the connection weight $W_{\alpha\beta}$, and the noise strength $\sigma_{\alpha\alpha}$. The detailed derivation of the conjugate force and the related response identities is summarized in Appendix~B.

\section{Explicit Response Formula}
\quad We write down the explicit response formulas corresponding to the steady-state linear response theory developed in the previous section. As an example, we consider the response of the equal-time covariance matrix $\mathbf{K}_0 \equiv \langle \delta \vec{x} \delta \vec{x}^\intercal \rangle$ with respect to the change in $\vec{p}$, and $\vec{p}$ can be $\vec{r}$, $\mathbf{W}$, and $\boldsymbol{\sigma}$. The observable becomes $B_{ij}(\vec{x}; \vec{p})=\delta x_i \delta x_j$. We also present the corresponding response functions for the different types of parameter perturbations. From the general formula (\ref{response_function2}), the response function with respect to $p_\alpha$ gives
\begin{equation}
    \chi_{ij\alpha} = \langle \partial_p (\delta x_i \delta x_j) \rangle_{ss, \vec{p}} + \langle \delta (\delta x_i \delta x_j) \delta \mathscr{F}_\alpha \rangle_{ss, \vec{p}}.
\end{equation}
The observable is written as $\delta x_i \delta x_j= [x_i - \langle x_i \rangle][x_j - \langle x_j \rangle]$. Then
\begin{equation}
    \left\langle \frac{\partial B_{ij}}{\partial p_\alpha} \right\rangle = -\frac{\partial \langle x_i \rangle}{\partial p_\alpha} \langle \delta x_j \rangle - \frac{\partial \left\langle x_j \right\rangle}{\partial p_\alpha} \langle \delta x_i \rangle.
\end{equation}
Since $\langle \delta x_i \rangle = 0$, the first term vanishes. The second term can be written as
\begin{eqnarray}
    \left\langle \delta (\delta x_i \delta x_j) \delta \mathscr{F}_\alpha \right\rangle_{ss, \vec{p}} &=& \left\langle (\delta x_i \delta x_j) \mathscr{F}_\alpha \right\rangle_{ss, \vec{p}} - \left\langle \delta x_i \delta x_j \right\rangle_{ss, \vec{p}} \left\langle \mathscr{F}_\alpha \right\rangle_{ss, \vec{p}}\nonumber\\
    &=& \langle (\delta x_i \delta x_j) \mathscr{F}_\alpha \rangle_{ss, \vec{p}} - (\mathbf{K}_0)_{ij} \langle \mathscr{F}_\alpha \rangle_{ss, \vec{p}}.
\end{eqnarray}
The explicit formulas of $\mathscr{F}_\alpha$ for three kinds of the network parameters are shown in Appendix~B.
\subsection{Response due to Change in the Intrinsic Node Parameter $r$}
\quad We first consider a perturbation of a single node perturbation $r_\alpha$, where $\alpha$ denotes the index of the perturbed node. One obtains the steady-state response of the equal-time covariance matrix
\begin{equation}
    \chi_{ij\alpha} = \langle \delta x_i \delta x_j \mathscr{F}_\alpha \rangle - (\mathbf{K}_0)_{ij} \langle \mathscr{F}_\alpha \rangle.
    \label{response_Ko_r}
\end{equation}
By using the dependence of $\mathbf{Q}$ on $r_\alpha$, the response function can be written as
\begin{eqnarray}
    \chi_{ij\alpha}^{(ss)} &=& \frac{\partial f'_\alpha}{\partial r_\alpha} \Big|_{\vec{X}}\left( \frac{\partial \mathbf{K}_0}{\partial Q_{\alpha\alpha}} \right)_{ij}-\frac{\partial f_\alpha}{\partial r_\alpha} \Big|_{\vec{X}} \sum_s Q^{-1}_{s\alpha} \Bigg[ f_s''\left( \frac{\partial \mathbf{K}_0}{\partial Q_{ss}} \right)_{ij}\nonumber\\
    &+& \sum_n W_{sn} \partial_1^2 h(X_s, X_n) \left( \frac{\partial \mathbf{K}_0}{\partial Q_{ss}} \right)_{ij}+\sum_\mu W_{\mu s} \partial_1 \partial_2 h(X_\mu, X_s) \left( \frac{\partial \mathbf{K}_0}{\partial Q_{\mu\mu}} \right)_{ij}\nonumber\\
    &+& \sum_\nu W_{s\nu} \partial_1\partial_2 h(X_s, X_\nu) \left( \frac{\partial \mathbf{K}_0}{\partial Q_{s\nu}} \right)_{ij} + \sum_\mu W_{\mu s} \partial^2_2 h(X_\mu, X_s) \left( \frac{\partial \mathbf{K}_0}{\partial Q_{\mu s}} \right)_{ij} \Bigg].
\end{eqnarray}
For the simple cases of $f_i(x_i) = r_i(a_i-x_i)$ with linear $h$, or Logistic $f_i(x_i) = r_i x_i (1-x_i)$ with synchronizing coupling and $h(X_i, X_j) = 0$, the above result simplifies significantly to 
\begin{equation}
    \chi_{ij\alpha}^{(ss)} = -\left(\frac{\partial \mathbf{K}_0}{\partial Q_{\alpha\alpha}} \right)_{ij}.
    \label{chiKodr_ss}
\end{equation}
This form shows that the response function of the covariance matrix is directly controlled by the diagonal element $Q_{\alpha\alpha}=-r_\alpha-\sum_j W_{\alpha j}$, namely by the local restoring parameter and the total incoming coupling associated with node $\alpha$. The derivative of $\mathbf{K}_0$ can be solved by differentiating the two sides of the Lyapunov equation (\ref{Lyapunov_eq}), which will be discussed later in the section. 

\subsection{Response due to Change in the Connection Weight $W$}
\quad We next consider the perturbation of a directed connection weight $W_{\alpha\beta}$. The corresponding response function of the covariance matrix is
\begin{equation}
    \chi_{ij(\alpha\beta)} = \langle \delta x_i \delta x_j \mathscr{F}_{\alpha\beta} \rangle - (\mathbf{K}_0)_{ij} \langle \mathscr{F}_{\alpha\beta} \rangle.
    \label{response_Ko_W}
\end{equation}
Again, by using the dependence of $\mathbf{Q}$ on $W_{\alpha\beta}$, the explicit formula is
\begin{eqnarray}
    \chi_{ij\alpha\beta}^{(ss)} &=& \partial_1 h(X_\alpha, X_\beta) \left( \frac{\partial \mathbf{K}_0}{\partial Q_{\alpha\alpha}} \right)_{ij} + \partial_2 h(X_\alpha, X_\beta) \left( \frac{\partial \mathbf{K}_0}{\partial Q_{\alpha\beta}} \right)_{ij}\nonumber\\
    &-& h(X_\alpha,X_\beta) \sum_s Q^{-1}_{s\alpha} \Bigg[ \Bigg( f''_s + \sum_n W_{sn} \partial^2_1 h(X_s, X_n) \Bigg) \left( \frac{\partial \mathbf{K}_0}{\partial Q_{ss}} \right)_{ij}\nonumber\\
    &+& \sum_\mu W_{\mu s} \partial_1\partial_2 h(X_\mu, X_s) \left( \frac{\partial \mathbf{K}_0}{\partial Q_{\mu\mu}} \right)_{ij} + \sum_\nu W_{s\nu} \partial_1\partial_2 h(X_s, X_\nu) \left( \frac{\partial \mathbf{K}_0}{\partial Q_{s\nu}} \right)_{ij}\nonumber\\
    &+& \sum_\mu W_{\mu s} \partial_2^2 h(X_\mu, X_s) \left( \frac{\partial \mathbf{K}_0}{\partial Q_{\mu s}} \right)_{ij} \Bigg].
\end{eqnarray}
For systems that can achieve a fully synchronized noise-free state with synchronizing coupling function $h(X_i, X_j)=0$, the formula reduces to
\begin{equation}
    \chi_{ij(\alpha\beta)} = \partial_1 h(X_\alpha,X_\beta) \left( \frac{\partial \mathbf{K}_0}{\partial Q_{\alpha\alpha}} \right)_{ij} + \partial_2 h(X_\alpha,X_\beta) \left( \frac{\partial \mathbf{K}_0}{\partial Q_{\alpha\beta}} \right)_{ij},
    \label{chiKodW_ss}
\end{equation}
which makes the role of the perturbed connection especially transparent. In this case, the response is determined by two contributions: the change in the local diagonal part $Q_{\alpha\alpha}$, and the change in the directed off-diagonal coupling $Q_{\alpha\beta}$. For an undirected network, since $W_{\alpha\beta}=W_{\beta\alpha}$, the total response is the sum of the two directed contributions associated with $W_{\alpha\beta}$ and $W_{\beta\alpha}$.

\subsection{Response due to Change in the Noise Strength $\sigma$}
\quad Finally, we consider the perturbation of a noise strength. In contrast to the previous two cases, a change in $\sigma_{\alpha\alpha}$ does not alter the deterministic force, and therefore does not shift the fixed point. The corresponding response function gives
\begin{equation}
    \chi_{ij(\alpha\alpha)} = \langle \delta x_i \delta x_j \mathscr{F}_{\alpha\alpha} \rangle - (\mathbf{K}_0)_{ij} \langle \mathscr{F}_{\alpha\alpha} \rangle = \left( \frac{\partial \mathbf{K}_0}{\partial \sigma_{\alpha\alpha}} \right)_{ij}.
\end{equation}
This result emphasizes that a perturbation of the noise strength affects the steady-state fluctuations directly, without changing the deterministic fixed point.

\subsection{Evaluation of $\partial \mathbf{K}_0/\partial Q_{\alpha\beta}$ and $\partial \mathbf{K}_0/\partial \sigma_{\alpha\alpha}$}
\quad The derivative of $\mathbf{K}_0$ with respect to $\mathbf{Q}$ or $\boldsymbol{\sigma}$ is determined from the steady-state Lyapunov equation due to the dependence of $\mathbf{K}_0(\mathbf{Q},\boldsymbol{\sigma})$. The derivatives of $\mathbf{K}_0$ with respect to the matrix elements of $\mathbf{Q}$ can be obtained by differentiating the Lyapunov equation
\begin{equation}
    \mathbf{Q} \frac{\partial \mathbf{K}_0}{\partial Q_{\alpha\beta}} + \frac{\partial \mathbf{K}_0}{\partial Q_{\alpha\beta}} \mathbf{Q}^\intercal = \mathbf{M}^{(\alpha,\beta)},\; (\mathbf{M}^{(\alpha,\beta)})_{ij} \equiv -(\mathbf{K}_0)_{i\beta} \delta_{j\alpha} - (\mathbf{K}_0)_{j\beta} \delta_{i\alpha}.
    \label{derivative_Ko}
\end{equation}
Also, the derivative of the inverse of the covariance matrix gives
\begin{equation}
    \mathbf{K}_0 \frac{\partial \mathbf{K}^{-1}_0}{\partial Q_{\alpha\beta}} \mathbf{K}_0 = - \frac{\partial \mathbf{K}_0}{\partial Q_{\alpha\beta}}\;\text{or}\;\mathbf{K}^{-1}_0 \frac{\partial \mathbf{K}_0}{\partial Q_{\alpha\beta}} \mathbf{K}^{-1}_0 = - \frac{\partial \mathbf{K}^{-1}_0}{\partial Q_{\alpha\beta}}.
    \label{inverse_partialKo}
\end{equation}
Higher-order derivatives can be obtained by the similar manner. For the derivative with respect to $\sigma_{\alpha\alpha}$, we have
\begin{equation}
    \mathbf{Q}\frac{\partial\mathbf{K}_0}{\partial\sigma_{\alpha\alpha}} + \frac{\partial\mathbf{K}_0}{\partial\sigma_{\alpha\alpha}}\mathbf{Q}^\intercal = -\mathbf{E}^{(\alpha\alpha)},\;\text{where}\;E^{(\alpha\alpha)}_{ij}=1\;\text{for}\;i=j=\alpha.
    \label{partialKo_sigma}
\end{equation}
Moreover, the above formulas are further simplified when the network satisfies the equilibrium condition $\mathbf{Q}\boldsymbol{\sigma}=\boldsymbol{\sigma}\mathbf{Q}^\intercal$. In that case, the covariance matrix satisfies $\mathbf{K}_0 = -\frac{1}{2} \mathbf{Q}^{-1} \boldsymbol{\sigma}$, and therefore
\begin{equation}
    \left( \frac{\partial \mathbf{K}_0}{\partial Q_{\alpha\beta}} \right)_{ij} = -Q^{-1}_{i\alpha} (\mathbf{K}_0)_{j\beta} = \frac{1}{2} Q^{-1}_{i\alpha} (\mathbf{Q}^{-1}\boldsymbol{\sigma})_{j\beta}.
\end{equation}
If the nodes share the same noise strength $\boldsymbol{\sigma}=\sigma \mathbf{I}$, the formula is further simplified
\begin{equation}
    \left( \frac{\partial \mathbf{K}_0}{\partial Q_{\alpha\beta}} \right)_{ij} = \frac{\sigma}{2} Q^{-1}_{i\alpha} Q^{-1}_{j\beta}.
\end{equation}
This equilibrium form is especially useful because the explicit dependence of $\mathbf{K}_0$ on $\mathbf{Q}$ is simple, so the response functions can be written directly in terms of $\mathbf{Q}^{-1}$. Then, the response function of $\mathbf{K}_0$ due to the node parameter change is simplified as
\begin{equation}
    \chi_{ij(\alpha)} = -\frac{\sigma}{2} Q^{-1}_{i\alpha} Q^{-1}_{j\alpha},
    \label{chiKodr_eq}
\end{equation}
and the response functions for the change in a connection weight gives
\begin{equation}
    \chi_{ij(\alpha\beta)} = -\frac{\sigma}{2} Q^{-1}_{i\alpha} Q^{-1}_{j\alpha}+\frac{\sigma}{2} Q^{-1}_{i\alpha} Q^{-1}_{j\beta},
    \label{chiKodW_eq}
\end{equation}
Notice that a noise perturbation usually breaks detailed balance. For a system originally at equilibrium, its perturbed state does not satisfy detailed balance. The equilibrium response formula of noise perturbations should be derived from Eq. (\ref{partialKo_sigma}). These results provide a simple reference form for comparison with the nonequilibrium steady-state response.

\section{Steady-State Response and Numerical validation}
\subsection{Simulation setup and observables}
\quad We perform Langevin simulations on synthetic networks to validate the theoretical response functions. 
Both undirected and directed networks are considered, with either uniform or heterogeneous noise. The network dynamics follow (\ref{noisy_network}), and the steady-state covariance matrix $\mathbf{K}_0 = \langle \delta \vec{x} \delta \vec{x}^\top \rangle$ is measured from long-time trajectories with time step $\Delta t = 0.0005$ and simulation period $T=10000$. To probe the system response, we introduce small perturbations to either node parameters $\delta r_\alpha$ or connection weights $\delta W_{\alpha\beta}$, where the perturbations of multiple nodes and edges are also considered to test the robustness, and these can be performed by small modifications to parameters. Disruptions of networks can be realized bu decreasing the connection weights of chosen edges to zero. For disruptions of nodes, they become isolated nodes where all edges incident on them are removed (connection weights of all edges connected with the chosen nodes become zero). After perturbation, the system relaxes to a new steady state, from which the updated covariance matrix is measured. The response is quantified by the difference between the perturbed and unperturbed steady states, $\Delta \mathbf{K}_0$, and compared with the theoretical prediction obtained from the unperturbed steady-state time series. We consider the undirected/directed ER random graphs with number of nodes $N=50, 100$ and the connection probability is about $0.1\sim0.2$. The intrinsic dynamics of each node is governed by the Logistic function $f_i=r_ix_i(1-x_i)$ and the interaction between nodes is the synchronizing coupling function $h(x_i,x_j)=x_j-x_i$. Under this condition, the Jacobian matrix $\mathbf{Q}$ is given by
\begin{equation}
    Q_{ii} = -r_i - \sum_{j\neq i} W_{ij},\; Q_{ij} = W_{ij},
\end{equation}
where the noise-free fixed points for each node is $1$. The network parameters are chosen to represent generic noisy network dynamics while allowing controlled comparisons between equilibrium and nonequilibrium steady states. From the detailed balance condition $\mathbf{Q}\boldsymbol{\sigma}=\boldsymbol{\sigma}\mathbf{Q}^\intercal$, an asymmetric $\mathbf{W}$ with a uniform diagonal $\boldsymbol{\sigma} = \sigma_0 \mathbf{I}$ and a symmtric $\mathbf{W}$ with a nonuniform diagonal $\boldsymbol{\sigma}$ always break detailed balance. The equilibrium case is the undirected graph with uniform noise strengths or the special choise $W_{ij} \sigma_{jj} = W_{ji} \sigma_{ii}$. As for the intrinsic node parameter, which does not affect detailed balance, are only chosen to ensure the existence of a stable fixed point. We consider the non-negative $r_i$ and $W_{ij}$ and the noise strenghs of all nodes are positive.\par
While the simulations on smaller networks clearly demonstrate the validity of the response formulas, 
it remains important to verify that the agreement is not restricted to finite-size or realization-specific effects. From a computational perspective, large networks are particularly relevant because the covariance matrix can be obtained efficiently from the Lyapunov equation, allowing the response functions to be evaluated without explicitly repeating perturbation simulations for every parameter change.Since the present framework is intended as a general response theory for noisy network dynamcis around fixed points, we also consider networks with a larger number of nodes and different classes of network topologies, including regular lattices, WS small-world networks, and scale-free networks. These network models represent distinct structural features, ranging from homogeneous local connectivity to highly heterogeneous degree distributions.

\subsection{Equilibrium response: undirected networks with uniform noise}
\quad We first consider undirected networks with uniform noise strengths, for which the steady state satisfies detailed balance. We perturb either node parameters $\delta r_\alpha$ or connection weights $\delta W_{\alpha\beta}$, and measure the resulting change in the steady-state covariance matrix $\mathbf{K}_0$. The theoretical prediction is obtained from the unperturbed steady-state statistics. Fig. \ref{chieqmlinear}a and b show the comparison between the simulation of the change in $\mathbf{K}_0$ and the theoretical response function with respect to the perturbation of a single node and edge, the response function is given as (\ref{chiKodr_eq}) and (\ref{chiKodW_eq}). The node parameters are uniformly distributed on $[0,5]$, and the connection weights are normally distributed $W_{ij} \sim \mathcal{N}(5,1)$. The noise strengths are identical among the nodes $\sigma_{ii}=0.01$. The data points collapse on the diagonal, which indicates the validation of the response function for noisy network dynamics at equilibrium for both node and edge perturbations. In Fig. \ref{chieqmlinear}a, the perturbed node ($\alpha$) is randomly chosen with three different positive perturbation values $dr_\alpha = 0.1, 0.5, 2$, denoted as circles, squares, and triangles. The node parameters determine the strength of the intrinsic forces on the nodes, so the increase of r stabilizes the nodes and reduces the effect of noise, and the elements in $\mathbf{K}_0$ decrease. The change in $(\mathbf{K}_0)_{ij}$ depends on the distance between the perturbed node $\alpha$ and the other nodes. The data point at the bottom left corner is $(\mathbf{K}_0)_{\alpha\alpha}$, which shows the perturbed node has the largest response. The change in $(\mathbf{K}_0)_{i\alpha}$ where node $i$ and the perturbed one has a direct connection $W_{i\alpha}$ and $W_{\alpha i} > 0$ is around $-5e-7$, and the other nodes that do not connect to $\alpha$ hardly respond to the perturbation.
\begin{figure}[H]
    \centering
    \subfigure[]{\includegraphics*[width=.45\columnwidth]{Figure/Steady/Undirected/dKdr.eps}}
    \subfigure[]{\includegraphics*[width=.45\columnwidth]{Figure/Steady/Undirected/dKdW.eps}}
    \caption{The two figures show the Langevin dynamics simulations for an undirected ER random network of $N=50$ nodes with linear force intrinsic node dynamics $f_i(x_i)=-r_i(x_i-1)$ and connection probability $p=0.2$ with a uniform noise matrix a normally distributed $\bf{W}$: (a) The figure shows the response function of $\bf{K}_0$ due to the change in $r_\alpha$ and $d r_\alpha = 0.1, 0.5, 2$ and is plotted against the simulation results. $\{r_i\}$ is normally distributed with the mean value $\overline{r}=5$. The point at the bottom left corner is the response of $({\bf K}_0)_{\alpha\alpha}$, which has the largest response. The points around $-5e-7$ are the nodes directly connected to $\alpha$, and those around 0 have no path to node $\alpha$. (b) The figure shows the response function of ${\bf K}_0$ due to the change in $W_{\alpha\beta}$, where the mean value $\overline{W} = 5$. The three points from bottom left to top right are the response of $({\bf K}_0)_{\beta\beta}$, $({\bf K}_0)_{\alpha\alpha}$, and $({\bf K}_0)_{\alpha\beta}$ respectively.}
    \label{chieqmlinear} 
\end{figure}
In Fig. \ref{chieqmlinear}b, the perturbed edge $(\alpha,\beta)$ is randomly chosen with three different positive perturbation values $dW_{\alpha\beta} = 0.2, 1, 5$, denoted as circles, squares, and triangles. Notice that for undrected networks, the $dW_{\beta\alpha}$ also changes to maintain the symmetric topology. The coupling strengths between any pair of nodes are characterized by the connection weights. The increase of the connection weight between node $\alpha$ and node $\beta$ enhances the correlation between the states of the two nodes. The data point at the top right corner is $(\mathbf{K}_0)_{\alpha\beta}$, and the auto-correlations $(\mathbf{K}_0)_{\alpha\alpha}$ and $(\mathbf{K}_0)_{\beta\beta}$ are at the left bottom corner, where the increasing coupling reduces the fluctuations of the two nodes. In the case of the Logistic force function, the response function of $\bf{K}_0$ is the same as that in the linear force function. As Fig. \ref{chieqmlogistic}a and \ref{chieqmlogistic}c show, the theoretical predictions match the simulation results when a parameter $r_\alpha$ or $W_{\alpha\beta}$ changes.
\begin{figure}[H]
    \centering
    \subfigure[]{\includegraphics*[width=.45\columnwidth]{Figure/Steady/Undirected/chiij1eqm.eps}}
    \subfigure[]{\includegraphics*[width=.45\columnwidth]{Figure/Steady/Undirected/LogisticKowun.eps}}
    \caption{Langevin dynamics simulations for an undirected ER random network of $N=50$ nodes with heterogeneous Logistic intrinsic node dynamics $f_i(x_i)=r_ix_i(1-x_i)$ and connection probability $p=0.2$, with a uniform diagonal noise matrix. (a) The response function of the fluctuation correlations. The change in the elements of the correlation functions upon a change of the node parameter of $\delta r_\alpha$ of the node $\alpha$ is measured by simulations. Four nodes are arbitrarily chosen with three different values of $\delta r_\alpha$. $\chi_{ij\alpha}$ is significantly larger for $i=j=\alpha$. (b) The response of ${\bf K}_0$ due to the change in $\bf W$ is large at $(i,j)=(\alpha,\alpha)$, $(\alpha,\beta)$, and $(\beta,\beta)$.}
    \label{chieqmlogistic} 
\end{figure}

\subsection{Nonequilibrium steady-state response}
\quad  In our framework, a NESS arises when detailed balance is broken, which can occur either through asymmetric interactions (directed networks) or heterogeneous noise distributions. In contrast to equilibrium systems, where the steady state is characterized by vanishing probability currents, NESS exhibit persistent irreversible fluxes. The key question we address is whether the steady-state linear response framework remains predictive under such conditions. To highlight the difference between equilibrium and nonequilibrium responses, we compare the exact response with the equilibrium-like simplified formula of $\mathbf{K}_0$, as given in (\ref{chiKodr_eq}) and (\ref{chiKodW_eq}). In Fig. \ref{chieqcompare}, a single node and a single directed edge are perturbed by $\delta r_\alpha, \delta W_{\alpha\beta}=0.5$ in a directed ER random network with a uniform noise matrix. The vertical axis is the equilibrium response function $\chi^{eq}_{ij\alpha(\beta)}$ and the horizontal one is the exact response function $\chi^{ss}_{ij\alpha(\beta)}$ solved from the Lyapunov equation. 
\begin{figure}[H]
    \centering
    \subfigure[]{\includegraphics*[width=.4\columnwidth]{Figure/Else/chieqnondr.eps}}
    \subfigure[]{\includegraphics*[width=.4\columnwidth]{Figure/Else/chieqnondW.eps}}
    \subfigure[]{\includegraphics*[width=.4\columnwidth]{Figure/Else/chieqdirdr.eps}}
    \subfigure[]{\includegraphics*[width=.4\columnwidth]{Figure/Else/chieqdirdW.eps}}    
    \caption{The comparison between the equilibrium and the NESS response functions for the covariance matrix due to the perturbations of a single node and an edge: When the detailed balance condition is violated, the response must instead be computed from the full Lyapunov response equation.}
    \label{chieqcompare} 
\end{figure}
Similar to the equilibrium case, we compare the simulation results with the theoretical response functions for the NESS conditions: a directed graph with a uniform diagonal noise matrix in Fig. \ref{chiKodir} and an undirected graph with a nonuniform diagonal noise matrix in Fig. \ref{chiKonon}. In both two cases, the perturbed node and the perturbed edge are randomly chosen, and the perturbed values are $\delta r_\alpha = 0.5, 2$ and $\delta W_{\alpha\beta} = 0.5, 2$. The distributions of the network parameters are the same as the equilibrium cases in Fig. \ref{chieqmlinear} and Fig. \ref{chieqmlogistic}. As the figures show, the response functions of the equal-time covariance are valid for the directed graph and the heterogenous noise strengths, where detailed balance is broken. Note that in the directed graph, we apply a single edge perturbation; however, in the undirected graph, when one edge $W_{\alpha\beta}$ is perturbed, the other edge $W_{\beta\alpha}$ with the opposite direction is also perturbed with the same amount to maintain the undirected topology after the perturbation. In Fig. \ref{chiKodir}b, the perturbed edge is from node $\beta$ to node $\alpha$, and thus only node $\alpha$ is affected. There is a data point $\delta (\mathbf{K}_0)_{\alpha\alpha} / \delta W_{\alpha\beta}$ at the bottom left corner. In Fig. \ref{chiKonon}b, both the reciprocal edges are perturbed with the same amount of the perturbation value, so node $\alpha$ and $\beta$ are perturbed. There are two data points at the bottom left corner, which indicates the response of $(\mathbf{K}_0)_{\alpha\alpha}$ and $(\mathbf{K}_0)_{\beta\beta}$ decrease due to the increase of the coupling between them.
\begin{figure}[H]
    \centering
    \subfigure[]{\includegraphics*[width=.45\columnwidth]{Figure/Steady/Directed/dKdr_dir.eps}}
    \subfigure[]{\includegraphics*[width=.43\columnwidth]{Figure/Steady/Directed/dKdW_dir.eps}}
    \caption{Comparison between Langevin simulations and the steady-state response prediction for a directed ER random network with a uniform diagonal noise matrix, where detailed balance is broken by asymmetric couplings. (a) Response of $\mathbf{K}_0$ to perturbations of a single node parameter with $\delta r_\alpha = 0.5, 2$. (b) Response of $\mathbf{K}_0$ to perturbations of a single directed edge weight with $\delta W_{\alpha\beta} = 0.5, 2$. The collapse of the data points onto the diagonal confirms the validity of the NESS response function obtained from the Lyapunov equation.}
    \label{chiKodir} 
\end{figure}

\begin{figure}[H]
    \centering
    \subfigure[]{\includegraphics*[width=.45\columnwidth]{Figure/Steady/Undirected/dKdr_non.eps}}
    \subfigure[]{\includegraphics*[width=.45\columnwidth]{Figure/Steady/Undirected/dKdW_non.eps}}
    \caption{Comparison between Langevin simulations and the steady-state response prediction for an undirected ER random network with a nonuniform diagonal noise matrix, where detailed balance is broken by heterogeneous noise strengths. (a) Response of $\mathbf{K}_0$ to perturbations of a single node parameter with $\delta r_\alpha = 0.5, 2$. (b) Response of $\mathbf{K}_0$ to perturbations of a single undirected edge, where both $W_{\alpha\beta}$ and $W_{\beta\alpha}$ are increased by the same amount $\delta W_{\alpha\beta} = 0.5, 2$ to preserve the symmetric topology. The agreement with the diagonal validates the NESS response function in the presence of heterogeneous noise.}
    \label{chiKonon} 
\end{figure}

\subsection{Multiple perturbations and removal of network components}
\quad After validating the linear-response theory for single-parameter perturbations, we next examine whether the response functions can be superposed to predict the effect of multiple simultaneous perturbations. We define the covariance change as
\begin{equation}
    \Delta (\mathbf{K}_0)_{ij} =
    (\mathbf{K}_0^{\rm after})_{ij}-(\mathbf{K}_0^{\rm before})_{ij}.
\end{equation}
For a set of perturbed nodes
\(\mathcal{P}_V=\{\alpha_1,\alpha_2,\ldots,\alpha_m\}\), the first-order prediction is
\begin{equation}
    \Delta (\mathbf{K}_0)_{ij}
    =
    \sum_{\alpha\in\mathcal{P}_V}
    \chi^{(r)}_{ij\alpha}\delta r_\alpha .
\end{equation}
Similarly, for a set of perturbed edges \(\mathcal{P}_E=\{(\alpha_1,\beta_1),(\alpha_2,\beta_2),\ldots, (\alpha_m,\beta_m)\}\), we use
\begin{equation}
    \Delta (\mathbf{K}_0)_{ij}
    =
    \sum_{(\alpha,\beta)\in\mathcal{P}_E}
    \chi^{(W)}_{ij(\alpha\beta)}\delta W_{\alpha\beta}.
\end{equation}
The perturbation amplitudes \(\delta r_\alpha\) and \(\delta W_{\alpha\beta}\) are allowed to be different for different perturbed nodes or edges. To test this superposition property, we perform Langevin simulations on undirected and directed ER networks with nonuniform noise sources.  In each realization, multiple nodes or edges are randomly selected and perturbed simultaneously. The covariance change measured from the simulations is then compared with the first-order prediction obtained by summing the corresponding single-parameter response functions.  Agreement with the diagonal line in the scatter plots indicates that the covariance response remains additive within the linear regime.\par
In Fig. \ref{MultipledKdr}, ten randomly selected nodes are perturbed with the same amount of perturbation values $\delta r_\alpha=0.5$, where the set of the perturbed nodes is denoted as $\mathcal{P}_V=\{ \alpha_1, \alpha_2, ..., \alpha_{10} \}$. For visual clarity, we separate the elements of the covariance matrix into diagonal and off-diagonal components. We also classify the elements according to whether their indices include nodes associated with the perturbation: diagonal elements are labeled as $\alpha\alpha$ or $ii,\;(i\notin \mathcal{P}_V)$, and off-diagonal elements as $i\alpha$ or $ij,\;(i,j \notin \mathcal{P}_V)$. The data points collapse onto the diagonal line, indicating quantitative agreement between the linear-response prediction and the numerical simulation for multiple node perturbations. In Fig. \ref{MultipledKdW}, twenty randomly selected undirected edges are perturbed with the same amount of values $\delta W_{\alpha\beta}=0.5$. The perturbed edges are represented as the set $\mathcal{P}_E = \{ (\alpha_k,\beta_k),\;k=1,2,...,20 \}$. Again, the data points are separated into the diagonal and offdiagonal elements. The elements are further classified as $\alpha\alpha$ or $ii,\;(i\neq \alpha)$ for the diagonal ones, and $i\alpha$ or $ij,\;(i,j\neq \alpha)$. Notice that the classification of edge perturbations is based on the receiving node $\alpha$, not on the source node $\beta$. The data points align along the diagonal, which proves the validation of the linear response framework on multiple edge perturbations.
\begin{figure}[H]
    \centering
    \subfigure[]{\includegraphics*[width=.45\columnwidth]{Figure/Multiple/MultipledKdrii_non.eps}}
    \subfigure[]{\includegraphics*[width=.45\columnwidth]{Figure/Multiple/MultipledKdrij_non.eps}}
    \caption{Comparison between Langevin simulations and the superposed linear-response prediction for simultaneous perturbations of ten node parameters in an ER random network with nonuniform noise. (a) Diagonal covariance changes $\Delta(\mathbf{K}_0)_{ii}$. (b) Off-diagonal covariance changes $\Delta(\mathbf{K}_0)_{ij}$. The grouping of data points distinguishes entries involving perturbed nodes from those that do not, and the collapse onto the diagonal confirms the additivity of the first-order response.}
    \label{MultipledKdr} 
\end{figure}
\begin{figure}[H]
    \centering
    \subfigure[]{\includegraphics*[width=.44\columnwidth]{Figure/Multiple/MultipledKdWii_non.eps}}
    \subfigure[]{\includegraphics*[width=.46\columnwidth]{Figure/Multiple/MultipledKdWij_non.eps}}
    \caption{Comparison between Langevin simulations and the superposed linear-response prediction for simultaneous perturbations of twenty undirected edges in an ER random network with nonuniform noise. (a) Diagonal covariance changes $\Delta(\mathbf{K}_0)_{ii}$. (b) Off-diagonal covariance changes $\Delta(\mathbf{K}_0)_{ij}$. The alignment of the data along the diagonal shows that the first-order responses to individual edge perturbations remain additive in the linear regime.}
    \label{MultipledKdW} 
\end{figure}
\begin{figure}[H]
    \centering
    \subfigure[]{\includegraphics*[width=.42\columnwidth]{Figure/Multiple/MultipledKdWii_remove.eps}}
    \subfigure[]{\includegraphics*[width=.46\columnwidth]{Figure/Multiple/MultipledKdWij_remove.eps}}
    \caption{Comparison between Langevin simulations and the first-order prediction for the removal of 40 directed edges in a directed ER random network with nonuniform noise. The theoretical prediction is obtained by summing the single-edge response functions with $\delta W_{\alpha\beta}=-W_{\alpha\beta}$ for all removed edges. (a) Diagonal covariance changes $\Delta(\mathbf{K}_0)_{ii}$. (b) Off-diagonal covariance changes $\Delta(\mathbf{K}_0)_{ij}$. The approximate collapse onto the diagonal shows that the linear response still captures the dominant effect of multiple edge removals.}
    \label{MultipledKdW_remove}
\end{figure}
We further consider removal events as a finite-amplitude stress test of the linear-response framework. Edge removal is implemented by setting \(W_{\alpha\beta}\to 0\), corresponding to \(\delta W_{\alpha\beta}=-W_{\alpha\beta}\). Node removal is implemented by isolating the selected node through the removal of all incoming and outgoing edges connected to it. Since these perturbations are not infinitesimal, the first-order prediction is not expected to be quantitatively exact. Instead, this comparison identifies the regime in which the linear response remains informative and the regime in which higher-order effects become important. Fig. \ref{MultipledKdW_remove} examines the removal of 40 directed edges in a directed ER network with nonuniform noise. The prediction is obtained by summing the single-edge response functions over all removed edges. The agreement with the diagonal line shows that, for this level of edge removal, the first-order response still captures the dominant covariance change, even in a nonequilibrium steady state generated by directed couplings and heterogeneous noise. Figure \ref{MultipledKdW_noderemove} considers the removal of a single node in an undirected network with nonuniform noise. Although only one node is removed, this operation simultaneously deletes all edges connected to that node and therefore represents a large structural perturbation. The visible deviations from the diagonal line are expected, especially for covariance elements involving the removed node or its neighbors. These deviations indicate the onset of higher-order response terms beyond the first-order response theory.
\begin{figure}[H]
    \centering
    \subfigure[]{\includegraphics*[width=.45\columnwidth]{Figure/Multiple/MultipledKdWii_noderemove.eps}}
    \subfigure[]{\includegraphics*[width=.45\columnwidth]{Figure/Multiple/MultipledKdWij_noderemove.eps}}
    \caption{Comparison between Langevin simulations and the first-order prediction for the removal of a single node in an undirected ER random network with nonuniform noise. The prediction is obtained by summing the responses of all incident edge removals associated with the deleted node. (a) Diagonal covariance changes $\Delta(\mathbf{K}_0)_{ii}$. (b) Off-diagonal covariance changes $\Delta(\mathbf{K}_0)_{ij}$. The visible deviations from the diagonal indicate the onset of higher-order effects for this large structural perturbation.}
    \label{MultipledKdW_noderemove}
\end{figure}
% \subsection{Large-scale networks and topology dependence}
% \quad To assess the scalability and topological applicability of the linear response framework, we next consider large-scale networks with $N > 1000$. In this regime, direct Langevin simulations become computationally expensive for estimating steady-state $\mathbf{K}_0$. We therefore validate the linear response predictions by comparing them with covariance changes obtained from direct solutions of the Lyapunov equation before and after parameter perturbations. This procedure tests whether the first-order response formula remains quantitatively accurate in large networks and whether its performance depends on network topology. We consider the synthetic networks such as Erdős-Rényi, Barabási-Albert, Watts-Strogatz, and ring lattices, which represent random, scale-free, small-world, and local nearest-neighbor connectivity, respectively. In Fig. \ref{Response_for_large_undirected_network}, multiple perturbations of nodes and edges are applied in the four types of undirected network topologies with $N=1000$. 40 randomly selected nodes are perturbed with the perturbation values on the intercal ($\delta r$ in $[-1, 1]$), and 100 undirected edges are perturbed with the perturbation values on the intercal ($\delta W$ in $[-1, 1]$). The mean values of the node parameters $r_i$ and connection weights $W_{ij}$ are fixed at $5$ and they are sampled from a uniform distribution on $[2.5, 7.5]$ and a Gaussian distribution $\mathcal{N}(5,0.5)$. The data points collapse on the diagonal for the comparison between the analytic solution of $\mathbf{K}_0$ from the Lyapunov equation and the response functions, these results show that the linear response framework and the superposition are also valid for the linearized noisy network dynamics with different topologies.

% \begin{figure}[H]
%     \centering
%     \subfigure[]{\includegraphics*[width=.45\columnwidth]{Figure/Large_network/ER_BA_undir_dr.eps}}
%     \subfigure[]{\includegraphics*[width=.45\columnwidth]{Figure/Large_network/WS_RING_undir_dr.eps}}
%     \subfigure[]{\includegraphics*[width=.47\columnwidth]{Figure/Large_network/ER_BA_undir_dW.eps}}
%     \subfigure[]{\includegraphics*[width=.45\columnwidth]{Figure/Large_network/WS_RING_undir_dW.eps}}
%     \caption{The }
%     \label{Response_for_large_undirected_network}
% \end{figure}

% After verifying the quantitative agreement between the theoretical response functions and numerical measurements, we next ask how the magnitude of the response is controlled by network structure. In particular, we focus on the local variance response $\Delta (\mathbf{K}_0)_{\alpha\alpha}$ and the pairwise covariance response $\Delta (\mathbf{K}_0)_{\alpha\beta}$. The former quantifies how the fluctuation amplitude of the perturbed or directly affected node changes, and therefore provides a measure of local sensitivity or robustness. The latter quantifies how the dynamical correlation between two nodes is modified by a node or edge perturbation, and therefore captures how local perturbations are transmitted through network interactions. Together, these two quantities separate the self-response of a node from the correlation response between coupled nodes.

\section{Time-Dependent Linear Response Theory in Noisy Network Dynamics}\label{sec:time-dependent}
\quad The steady-state response developed above characterizes the difference between two nearby stationary states before and after a small parameter perturbation. In many situations, however, a perturbation is applied through a finite-time protocol rather than instantaneously. We therefore consider parameter protocols that drive the system from one steady state to another. Starting from a system that initially fluctuates near a fixed point, we derive a time-dependent response theory that predicts the average change of observables as a function of time. In particular, we derive the responses of the mean state, deterministic force, and equal-time covariance to time-dependent perturbations of the network parameters. We then validate the theory numerically for linear network dynamics in both an equilibrium reference network and a nonequilibrium steady-state network.
\subsection{General Formula for Time-Dependent Linear Response Function}
\quad We consider the response induced by a time-dependent perturbation of the network parameters, $\vec{p}\rightarrow\vec{p}+\delta \vec{p}(t)$, and assume that the system is in the unperturbed steady state for $t\le 0$. The time-dependent linear-response function of an observable $\mathbf{B}(t)$ is defined as
\begin{equation}
    \langle \mathbf{B}(t) \rangle_{\vec{p}+\delta \vec{p}(t)} = \langle \mathbf{B}(0) \rangle_{ss,\vec{p}} + \int_0^t dt' \boldsymbol{\chi}(t-t') \delta \dot{\vec{p}}(t').
    \label{time-dependent_response}
\end{equation}
The conjugate force $\vec{\mathscr{F}}(t)$ is defined in the same way as in Eq. (\ref{conjugate_force}). The time-dependent linear-response function is then given by
\begin{eqnarray}
    \boldsymbol{\chi}(t) &=&\langle \nabla_p{\mathbf{B}} \rangle_{ss,{\vec p}} + \int_0^t dt'\langle {\mathbf {B}}(t-t'){\dot{\vv {\mathscr F}}}^\intercal\rangle_{ss,{\vec p}} \nonumber\\
    &=& \boldsymbol{\chi} ^{(ss)}- \langle \delta{\mathbf B}(t)\delta{\vv {\mathscr F}}^\intercal(0)\rangle_{ss,{\vec p}}=\boldsymbol{\chi} ^{(ss)}- \langle \delta{\mathbf B}(t){\vv {\mathscr F}}^\intercal(0)\rangle_{ss,{\vec p}}.
    \label{chineqm}
\end{eqnarray}
The corresponding time-dependent response of $\mathbf{B}(t)$ is
\begin{equation}
    \langle \mathbf{B}(t)\rangle_{{\vec p}+\delta{\vec p}(t)}= \langle \mathbf{ B}(0)\rangle_{ss,{\vec p}}+\boldsymbol{\chi}^{(ss)}\delta {\vec p}(t)-\int_0^t dt'\langle \delta\mathbf{ B}(t-t'){\vv {\mathscr F}}^\intercal(0)\rangle_{ss,{\vec p}} \delta{\dot {\vec p}}(t').
    \label{Bnoneqm}
\end{equation}
For the linearized noisy dynamics, both the time-lagged covariance $\mathbf{K}_\tau$ and the mixed correlation with the conjugate force are governed by $e^{\tau \mathbf{Q}}$. Using the time-dependent conjugate-force relation derived in Appendix~B, one obtains
\begin{equation}
    \langle \delta x_i(\tau)\mathscr{F}_\alpha(0)\rangle = \sum_\ell \frac{\partial X_\ell}{\partial p_\alpha}
    (e^{\tau \mathbf{Q}})_{i\ell}.
    \label{time_mixed_xF}
\end{equation}

This result shows that, within the linearized theory, the time-dependent response is determined by two ingredients: the parameter-induced shift of the noise-free fixed point and the propagator $e^{\tau\mathbf{Q}}$, which governs how the perturbation is transmitted through the network in time.
\subsection{Response Functions in the Linearized-Dynamics Framework}
\quad We now specialize the general time-dependent response formulas to the linearized dynamics around the stable fixed point. The response functions below follow from Eqs. (\ref{chineqm}) and (\ref{time_mixed_xF}) within this linearized Gaussian framework. Eqs. (\ref{chinsslinear})--(\ref{chitKolinear}) assume the local linearization around the fixed point.\par
For the mean state, combining Eq. (\ref{chineqm}) with Eq. (\ref{time_mixed_xF}) gives
\begin{equation}
    \chi_{i\alpha}^{(x)}(t)=\left(e^{t{\bf Q}}{\bf Q}^{-1} \nabla_p{\vec { F}}|_{\vec X} \right)_{i\alpha}-\left({\bf Q}^{-1} \nabla_p{\vec { F}} |_{\vec X}\right)_{i\alpha}. \label{chinsslinear}
\end{equation}
For the deterministic force, the linearized fluctuation satisfies $\delta \vec{F}(t)=\mathbf{Q}\,\delta \vec{x}(t)$ to leading order. The corresponding time-dependent response function is therefore
\begin{equation}
    \chi_{i\alpha}^{(F)}(t)=\left(e^{t{\bf Q}}\nabla_p{\vec { F}}|_{\vec X} \right)_{i\alpha}.
    \label{chitforcelinear}
\end{equation}
For the equal-time covariance matrix $\mathbf{K}(t)\equiv \langle \delta \vec{x}(t) \delta \vec{x}^\intercal(t) \rangle$, its time-dependent response function gives
\begin{equation}
    \chi_{ij\alpha}^{(K)}(t) = \chi_{ij\alpha}^{(ss)} - \sum_{\ell,m} \left((e^{t\mathbf{Q}})_{i\ell}\chi_{\ell m\alpha}^{(ss)}(e^{t\mathbf{Q}^\intercal})_{mj}\right),
    \label{chitKolinear}
\end{equation}
where $\chi_{\ell m\alpha}^{(ss)}$ is the steady-state response function derived in Section IV. These expressions provide the explicit time-dependent response kernels within the linearized framework. We next specify the simulation setting and then evaluate them further for the particular linear force and diffusive coupling used in the numerical tests.
\subsection{Simulation Setup and Validation Strategy}
To validate the explicit time-dependent formulas derived for linearized network dynamics, we consider the two four-node networks shown in Fig. \ref{graphN4}: an undirected network satisfying the equilibrium condition and a directed network realizing a nonequilibrium steady state. In both cases, the node parameters and nonzero coupling weights are fixed at $r_i=5$ and $W_{ij}=5$, respectively. The intrinsic dynamics is chosen as
\begin{equation}
    f_i(x_i)=r_i(a_i-x_i),
\end{equation}
with intrinsic target states $\vec{a}=(2,4,6,8)$, and the coupling function is diffusive,
\begin{equation}
    h(x_i,x_j)=x_j-x_i .
\end{equation}
For the undirected network, all nodes have identical noise strengths, $\sigma_{ii}=0.01$, so the system serves as an equilibrium reference case. In contrast, the directed network contains four unidirectional edges, thereby producing a nonequilibrium steady state.\par
We perturb either the node parameter $r_1(t)$ or a single coupling weight $W_{12}(t)$, corresponding to the edge from node 2 to node 1 under our convention. For the undirected network, the same perturbation is applied to the reciprocal edge between nodes 1 and 2. The time-dependent perturbation protocols are given in Eqs. \eqref{dra(t)cos} and \eqref{dWab(t)cos}, with amplitude $b=0.5$ and protocol duration $\tau=10$. The same protocol parameters are used for all time-dependent comparisons shown below. In the following subsections, we evaluate the corresponding explicit linear-response formulas for single-node and single-edge perturbations and compare them with direct Langevin simulations.

The observables used for validation are the mean state $\langle x_i(t)\rangle$, 
the mean force $\langle F_i(t)\rangle$, and the equal-time covariance matrix
\begin{equation}
    K_{ij}(t,t) = \left\langle \delta x_i(t)\delta x_j(t)\right\rangle .
\end{equation}
All ensemble averages are computed from $10^5$ independent trajectories over the time interval $0\leq t\leq 20$. Before the perturbation is applied at $t=0$, the system is first allowed to relax to its unperturbed steady state. In the time-series comparisons below, dashed curves denote direct Langevin simulations and solid curves denote the linear-response prediction. Since showing all components would be visually redundant for this four-node example, we display the representative quantities $\langle x_1(t)\rangle$, $\langle x_2(t)\rangle$, $\langle F_1(t)\rangle$, $\langle F_2(t)\rangle$, $K_{11}(t,t)$, and $K_{12}(t,t)$. The remaining components show similarly good agreement, so only representative responses are plotted.

\begin{figure}[H]
    \centering
    \subfigure[]{\includegraphics*[width=.45\columnwidth]{Figure/Time_dependent/undirected_graph_N4.eps}}
    \subfigure[]{\includegraphics*[width=.45\columnwidth]{Figure/Time_dependent/directed_graph_N4.eps}} 
    \caption{Four-node networks used to validate the time-dependent response formulas. (a) Undirected cycle with uniform noise, which serves as an equilibrium reference case. (b) Directed cycle with heterogeneous noise, which realizes a nonequilibrium steady state. In the numerical tests, we perturb either node $1$ through $\delta r_1(t)$ or the directed edge from node $2$ to node $1$ through $\delta W_{12}(t)$; for the undirected network, the reciprocal edge is changed by the same amount to preserve the symmetric topology.}
    \label{graphN4} 
\end{figure}

\subsection{Time-dependent Response due to the time-dependent perturbation in a node parameter}
\quad We first consider a time-dependent perturbation of a single node parameter $\delta r_\alpha(t)$. For $f_i(x_i;r_i) = r_i(a_i-x_i)$ and $h(x_i,x_j)=x_j-x_i$, the time-dependent response of $x_i$ is
\begin{equation}
    \langle x_i(t)\rangle_{r+\delta r_\alpha(t)}=X_i+ (a_\alpha-X_\alpha)\bigg[\int_0^tdt'\left(e^{(t-t'){\bf Q}}{\bf Q}^{-1}  \right)_{i\alpha}\delta{\dot r}_\alpha(t') -Q^{-1}_{i\alpha} \delta r_\alpha(t)\bigg].\label{xi(t)drlinear}
\end{equation}
In the numerical tests, we use the cosine protocol
\begin{equation}
     \delta r_\alpha(t)=\begin{cases} \frac{b}{2}(1-\cos\frac{\pi t}{\tau}), & 0\leq t\leq\tau\\
     b &t>\tau
     \end{cases},\label{dra(t)cos}
\end{equation}
As $t\to\infty$, since the eigenvalues of $\mathbf{Q}$ are negative, $\langle x_i(t)\rangle$ approaches the new noise-free fixed point. Using Eq. (\ref{chitforcelinear}), the mean-force response is
\begin{equation}
    \langle F_i(t)\rangle_{r+\delta r_\alpha(t)}
    =(a_\alpha-X_\alpha)\int_0^t dt'
    \left(e^{(t-t'){\bf Q}}\right)_{i\alpha}\delta{\dot r}_\alpha(t').
    \label{Fi(t)drlinear}
\end{equation}
The ensemble averages of $\vec{x}(t)$ and $\vec{F}(t)$ do not depend on the noise strengths and are determined solely by ${\bf Q}$. The equal-time covariance response is
\begin{equation}
    K_{ij}(t,t)=({\bf K_0})_{ij}
    -\left( \frac{\partial {\bf K}_0}{\partial Q_{\alpha\alpha}} \right)_{ij}\delta r_\alpha(t)
    +\int_0^t dt' \left( e^{(t-t'){\bf Q}}\frac{\partial {\bf K}_0}{\partial Q_{\alpha\alpha}} e^{(t-t'){\bf Q}^\intercal} \right)_{ij}\delta \dot{r}_\alpha(t').
    \label{Kij(t)drlinear}
\end{equation}
where $\mathbf{K}_0$ denotes the steady-state covariance for $t\leq 0$. For the cosine protocol (\ref{dra(t)cos}), the convolution integrals in Eqs. (\ref{xi(t)drlinear})--(\ref{Kij(t)drlinear}) can be analyzed by spectral decomposition of ${\bf Q}$. Appendix C gives the resulting closed-form expressions for $\langle x_i(t)\rangle$ and $\langle F_i(t)\rangle$, together with the mode-expanded covariance response written in terms of the protocol integral $I_{m\ell}(t)$.
\begin{figure}[H]
    \centering
    \subfigure[]{\includegraphics*[width=.45\columnwidth]{Figure/Time_dependent/nonsteady_response_xunddr.eps}}
    \subfigure[]{\includegraphics*[width=.47\columnwidth]{Figure/Time_dependent/nonsteady_response_funddr.eps}} 
    \subfigure[]{\includegraphics*[width=.45\columnwidth]{Figure/Time_dependent/nonsteady_response_Kounddr.eps}}
    \caption{Comparison between Langevin simulations (dashed) and the time-dependent linear-response prediction (solid) for a node perturbation $\delta r_1(t)$ in the four-node undirected network of Fig. \ref{graphN4}a. (a) Mean-state response $\langle x_i(t)\rangle-\langle x_i(0)\rangle$ for the components $i=1,2$. (b) Mean-force response $\langle F_i(t)\rangle$ for $i=1,2$. (c) Equal-time covariance response for the components $K_{11}(t,t)$ and $K_{12}(t,t)$. The agreement over both the driving stage and the post-driving relaxation verifies the time-dependent response theory in the equilibrium case.}
    \label{Time-dependent_undirected_dr} 
\end{figure}
Figures \ref{Time-dependent_undirected_dr} and \ref{Time-dependent_directed_dr} compare these predictions with Langevin simulations for a perturbation of the node parameter $r_1(t)$. The curves resolve both stages of the response: the driven interval $0\le t\le \tau$, during which the observables follow the smooth protocol, and the subsequent relaxation for $t>\tau$, during which the system approaches the new steady state. Agreement for $\langle x_i(t)\rangle$ and $\langle F_i(t)\rangle$ verifies the predicted mean relaxation dynamics. Agreement for $K_{ij}(t,t)$ shows that the theory also captures the time evolution of fluctuations and correlations.
\begin{figure}[H]
    \centering
    \subfigure[]{\includegraphics*[width=.45\columnwidth]{Figure/Time_dependent/nonsteady_response_xdirdr.eps}}
    \subfigure[]{\includegraphics*[width=.47\columnwidth]{Figure/Time_dependent/nonsteady_response_fdirdr.eps}} 
    \subfigure[]{\includegraphics*[width=.45\columnwidth]{Figure/Time_dependent/nonsteady_response_Kodirdr.eps}}
    \caption{Same comparison as in Fig. \ref{Time-dependent_undirected_dr}, but for the directed network in Fig. \ref{graphN4}b, which is in a nonequilibrium steady state. The observables are again (a) $\langle x_1(t)\rangle-\langle x_1(0)\rangle$ and $\langle x_2(t)\rangle-\langle x_2(0)\rangle$, (b) $\langle F_1(t)\rangle$ and $\langle F_2(t)\rangle$, and (c) $K_{11}(t,t)$ and $K_{12}(t,t)$. The close overlap between theory and simulation shows that the time-dependent response formulas remain valid even when detailed balance is broken.}
    \label{Time-dependent_directed_dr} 
\end{figure}
Increasing $r_1$ strengthens the local restoring force on node 1 toward its intrinsic target $a_1$. Since the unperturbed state satisfies $X_1>a_1$ because of the couplings, the perturbation drives $\langle x_1(t)\rangle$ downward and gives a negative response in $\langle F_1(t)\rangle$. The stronger local relaxation also suppresses fluctuations around node 1, leading to a decrease in $K_{11}(t,t)$. In the undirected case, the perturbation is transmitted more directly to node 2, so both the coupled response and the covariance $K_{12}(t,t)$ decrease. In the directed case, node 2 is not directly affected by node 1, and its response can therefore differ in amplitude and transient shape.
\subsection{Time-dependent Response due to the time-dependent perturbation in a connection weight}
\quad We next consider a time-dependent perturbation of a single connection weight. The derivation parallels the node-perturbation case, so we state the corresponding formulas for the same linear force and diffusive coupling used in the simulations and then compare them with direct simulations.
For this specific network model, the relaxation response of the mean position is
\begin{equation}
    \langle x_i(t)\rangle_{{\bf W}+\delta W_{\alpha\beta}(t)}
    =X_i+ (X_\beta-X_\alpha)\bigg[\int_0^t dt'\left(e^{(t-t'){\bf Q}}{\bf Q}^{-1} \right)_{i\alpha}\delta{\dot W}_{\alpha\beta}(t') -Q^{-1}_{i\alpha} \delta W_{\alpha\beta}(t)\bigg].
    \label{xi(t)dWlinear}
\end{equation}
We again use the cosine protocol on the connection between node $\alpha$ and node $\beta$:
\begin{equation}
     \delta W_{\alpha\beta}(t)=\begin{cases} \frac{b}{2}(1-\cos\frac{\pi t}{\tau}), & 0\leq t\leq\tau\\
     b &t>\tau
     \end{cases},\label{dWab(t)cos}
\end{equation}
The mean-force response is
\begin{equation}
    \langle F_i(t)\rangle_{{\bf W}+\delta W_{\alpha\beta}(t)}
    =(X_\beta-X_\alpha)\int_0^t dt'
    \left(e^{(t-t'){\bf Q}}\right)_{i\alpha}\delta{\dot W}_{\alpha\beta}(t').
    \label{Fi(t)dWlinear}
\end{equation}
The equal-time covariance response is
\begin{eqnarray}
    K_{ij}(t,t)&=&({\bf K_0})_{ij}
    +\left[- \frac{\partial {\bf K}_0}{\partial Q_{\alpha\alpha}}+\frac{\partial {\bf K}_0}{\partial Q_{\alpha\beta}} \right]_{ij}\delta W_{\alpha\beta}(t)\nonumber\\
    &+&\int_0^t dt' \left[ e^{(t-t'){\bf Q}}\left(-\frac{\partial {\bf K}_0}{\partial Q_{\alpha\alpha}} +\frac{\partial {\bf K}_0}{\partial Q_{\alpha\beta}}\right) e^{(t-t'){\bf Q}^\intercal} \right]_{ij} \delta \dot{W}_{\alpha\beta}(t').
    \label{Kij(t)dWlinear}
\end{eqnarray}
For the cosine protocol (\ref{dWab(t)cos}), the convolution integrals in Eqs. (\ref{xi(t)dWlinear})--(\ref{Kij(t)dWlinear}) can be evaluated explicitly by the same spectral decomposition of ${\bf Q}$ used in Appendix C. Figures \ref{Time-dependent_undirected_dW} and \ref{Time-dependent_directed_dW} show the corresponding validation for a time-dependent edge perturbation. The response is driven by the contrast $X_\beta-X_\alpha$ across the perturbed link, so the transient signals directly probe how a local change in coupling propagates through the network and modifies both the deterministic relaxation and the covariance dynamics.\par
Increasing $W_{12}$ strengthens the pull of node 2 on node 1. Since the unperturbed state satisfies $X_2>X_1$, the perturbation drives $\langle x_1(t)\rangle$ upward and gives a positive response in $\langle F_1(t)\rangle$. In the undirected case, the same perturbation also acts on $W_{21}$, reduces the state difference between the two nodes, and enhances their correlation, so $K_{12}(t,t)$ increases while $K_{11}(t,t)$ decreases. In the directed case, the response of node 2 is not directly affected by $\delta W_{12}$, so its amplitude is smaller and its time dependence is set by the network propagation.
\begin{figure}[H]
    \centering
    \subfigure[]{\includegraphics*[width=.45\columnwidth]{Figure/Time_dependent/nonsteady_response_xunddW.eps}}
    \subfigure[]{\includegraphics*[width=.47\columnwidth]{Figure/Time_dependent/nonsteady_response_funddW.eps}} 
    \subfigure[]{\includegraphics*[width=.45\columnwidth]{Figure/Time_dependent/nonsteady_response_KounddW.eps}}
    \caption{Comparison between Langevin simulations (dashed) and the time-dependent linear-response prediction (solid) for an edge perturbation in the four-node undirected network. The coupling on the link between nodes $1$ and $2$ is varied according to $\delta W_{12}(t)=\delta W_{21}(t)$. (a) Mean-state response for $\langle x_1(t)\rangle-\langle x_1(0)\rangle$ and $\langle x_2(t)\rangle-\langle x_2(0)\rangle$. (b) Mean-force response for $\langle F_1(t)\rangle$ and $\langle F_2(t)\rangle$. (c) Equal-time covariance response for $K_{11}(t,t)$ and $K_{12}(t,t)$. The agreement confirms the predicted transient and long-time response for a time-dependent perturbation of an undirected coupling.}
    \label{Time-dependent_undirected_dW} 
\end{figure}
\begin{figure}[H]
    \centering
    \subfigure[]{\includegraphics*[width=.45\columnwidth]{Figure/Time_dependent/nonsteady_response_xdirdW.eps}}
    \subfigure[]{\includegraphics*[width=.47\columnwidth]{Figure/Time_dependent/nonsteady_response_fdirdW.eps}} 
    \subfigure[]{\includegraphics*[width=.45\columnwidth]{Figure/Time_dependent/nonsteady_response_KodirdW.eps}}
    \caption{Same comparison as in Fig. \ref{Time-dependent_undirected_dW}, but for a directed perturbation $\delta W_{12}(t)$ in the nonequilibrium network of Fig. \ref{graphN4}b. Panels (a)-(c) show components of the mean state, mean force, and equal-time covariance, respectively. The close agreement between simulation and prediction demonstrates that the time-dependent response theory also captures the propagation of directed coupling perturbations in a nonequilibrium steady state.}
    \label{Time-dependent_directed_dW} 
\end{figure}
\subsection{Time-dependent Response due to the time-dependent perturbation in a noise strength}
\quad We next consider a time-dependent perturbation of the diagonal noise strength, $\delta \sigma_{\alpha\alpha}(t)$. Unlike node and edge perturbations, a perturbation in $\sigma_{\alpha\alpha}$ does not modify the deterministic drift matrix ${\bf Q}$. Therefore, for linear dynamics with additive zero-mean noise, the ensemble averages of the mean state and deterministic force remain unchanged to first order. In other words, a perturbation of the noise strength affects only the fluctuation sector at linear order, while the nontrivial response appears in the equal-time covariance matrix. Using Eq.~(\ref{Bnoneqm}), the time-dependent covariance response is
\begin{equation}
    {\bf K}(t,t)
    =
    {\bf K}_0
    +
    \frac{\partial {\bf K}_0}{\partial \sigma_{\alpha\alpha}}
    \,\delta \sigma_{\alpha\alpha}(t)
    -
    \int_0^t dt'\,
    e^{(t-t'){\bf Q}}
    \frac{\partial {\bf K}_0}{\partial \sigma_{\alpha\alpha}}
    e^{(t-t'){\bf Q}^\intercal}
    \,\delta \dot{\sigma}_{\alpha\alpha}(t').
    \label{Ksigma_time_response}
\end{equation}
This expression has the same structure as in the node- and edge-perturbation cases, except that the kernel is now determined by $\partial {\bf K}_0/\partial \sigma_{\alpha\alpha}$ rather than by derivatives with respect to entries of ${\bf Q}$. With the spectral decomposition ${\bf Q}=\sum_m \lambda_m {\vec u_m}{\vec v_m}^\intercal$, the covariance response can be written as
\begin{equation}
    K_{ij}(t,t)
    =
    (K_0)_{ij}
    +
    \left(
    \frac{\partial {\bf K}_0}{\partial \sigma_{\alpha\alpha}}
    \right)_{ij}
    \delta \sigma_{\alpha\alpha}(t)
    -
    \sum_{m,\ell}
    u_{im}u_{j\ell}
    \left(
    {\vec v_m}^\intercal
    \frac{\partial {\bf K}_0}{\partial \sigma_{\alpha\alpha}}
    {\vec v_\ell}
    \right)
    I_{m\ell}(t),
    \label{Ksigma_mode_expansion}
\end{equation}
where the protocol integral $I_{m\ell}(t)$ is the same as that defined in Eq.~(\ref{Iml}), with $\delta r_\alpha(t)$ replaced by $\delta \sigma_{\alpha\alpha}(t)$. For the cosine protocol used in the numerical test, we apply a time-dependent perturbation to $\sigma_{11}$ as
\begin{equation}
    \delta \sigma_{\alpha\alpha}(t)=\begin{cases} \frac{b}{2}(1-\cos\frac{\pi t}{\tau}), & 0\leq t\leq\tau \\
              b                                      & t>\tau
    \end{cases}.\label{dsaa(t)cos}
\end{equation}
Figure \ref{Time-dependent_Ko_ds} compares the predicted and simulated time-dependent responses of $K_{11}(t,t)$ and $K_{12}(t,t)$ for the two four-node graphs in Fig. \ref{graphN4}. Since the perturbation acts directly on the noise strength at node $1$, the most pronounced effect appears in the local variance $K_{11}(t,t)$. The cross-covariance $K_{12}(t,t)$ responds more indirectly through the propagation of fluctuations across the network. In the directed network of Fig. \ref{graphN4}b, node $2$ is not directly driven by node $1$, so the response of $K_{12}(t,t)$ is weaker than in the undirected case. In contrast, the bidirectional couplings in Fig. \ref{graphN4}a transmit the noise perturbation more efficiently, leading to a stronger covariance response.
\begin{figure}[H]
    \centering
    \subfigure[]{\includegraphics*[width=.45\columnwidth]{Figure/Time_dependent/nonsteady_response_Ko11ds.eps}}
    \subfigure[]{\includegraphics*[width=.45\columnwidth]{Figure/Time_dependent/nonsteady_response_Ko12ds.eps}}
    \caption{Comparison between Langevin simulations (dashed) and the time-dependent linear-response prediction (solid) for a noise-strength perturbation $\delta \sigma_{11}(t)$ in the two four-node networks of Fig. \ref{graphN4}. Panel (a) shows the equal-time variance $K_{11}(t,t)$ and panel (b) shows the equal-time covariance $K_{12}(t,t)$. In each panel, the responses for the undirected equilibrium network [Fig. \ref{graphN4}a] and the directed nonequilibrium network [Fig. \ref{graphN4}b] are plotted together. The close agreement confirms that a time-dependent perturbation of the local noise strength is captured accurately by the covariance response formula, while the weaker $K_{12}(t,t)$ signal in the directed case reflects the reduced propagation of fluctuations from node $1$ to node $2$.}
    \label{Time-dependent_Ko_ds}
\end{figure}

\section{Conclusion and Outlook}
\quad In this work, we developed a computable linear-response framework for noisy network dynamics near stable fixed points. Linearizing the stochastic dynamics around the noise-free fixed point reduces the problem to a Jacobian matrix ${\bf Q}$ and a steady-state covariance matrix ${\bf K}_0$ satisfying a Lyapunov equation. Within this structure, we derived first-order response formulas for perturbations of intrinsic node parameters, connection weights, and noise strengths. The main result is an operational theory for the equal-time covariance response that can be evaluated directly from the unperturbed linearized network and applied to both equilibrium and nonequilibrium steady states with additive Gaussian white noise.\par
The results show that linear response remains predictive for fluctuation observables even when detailed balance is broken. In particular, the covariance response gives a direct way to quantify how local perturbations spread through network-wide fluctuations and correlations. Our Langevin simulations confirmed the steady-state predictions in undirected equilibrium networks and in nonequilibrium networks generated by asymmetric couplings or heterogeneous noise strengths. Tests of multiple simultaneous perturbations further showed that the first-order responses remain additive within the linear regime. The removal examples showed where the first-order approximation begins to break down and thus point to the need for a nonlinear response theory for large perturbations and node removal. Since local perturbations propagate through correlations across the whole network, it will also be important to clarify more systematically how the response depends on topology.\par
We also extended the framework to time-dependent perturbations and obtained explicit time-dependent response formulas for the mean state, deterministic force, and equal-time covariance. This extension shows that the same formalism can describe not only the difference between two steady states, but also the full relaxation process induced by a driving protocol. In this picture, a parameter perturbation drives the system from one steady state to a new one, and the transient response describes how the system moves between them. This viewpoint suggests a natural connection to nonequilibrium steady-state thermodynamics and to the Hatano-Sasa description of transitions between steady states \cite{10.1143/PTPS.130.29,PhysRevLett.86.3463,doi:10.1073/pnas.0406405101}. Another natural extension is to relax the assumption of additive white noise with diagonal covariance. In many realistic systems, fluctuations are temporally correlated or shared across components, leading to colored noise \cite{PhysRevLett.84.3017, PhysRevResearch.5.033208, qrvl-mng6} and noise matrices with off-diagonal elements. Extending the framework to such cases would broaden its applicability and clarify how temporal memory and cross-component noise correlations reshape the response. For nonequilibrium steady states, another direction is to include the energetics of stochastic trajectories \cite{10.1143/PTPS.130.17} and connect the present covariance-based formalism to excess dissipation, housekeeping dissipation, and entropy production \cite{Seifert_2010, Seifert_2012, Mehl_2010}. Since the present theory already gives explicit conjugate-force and time-dependent response kernels in terms of ${\bf Q}$, ${\bf K}_0$, and the perturbation protocol, it should be possible to compute the linear response of these energetic observables directly for noisy networks. Finally, because the current theory is built around fluctuations near stable fixed points, an important next step is to generalize it to oscillatory regimes such as stable limit cycles \cite{PhysRevE.108.064114}.

\appendix
\section{Noisy network reconstruction formulas}
\quad Starting from Eq. (\ref{linearized_dynamics}), we have
\begin{equation}
\Delta{\dot{\vec x}} ={\bf Q}\Delta{\vec x} +{\vec \eta}(t)+{\cal O}(\Delta x ^2)\label{deltax}.
\end{equation}
For stable noise-free fixed points, the real parts of all eigenvalues of ${\bf Q}$ cannot be positive.
${\bf Q}$ is time-independent, and  the linearized motion in (\ref{deltax}) can be solved to give
\begin{equation}
    \Delta{\vec x}(t)=e^{ t{\bf Q}}\Delta{\vec x}(0)+\int^t_0 dt' e^{ (t-t'){\bf Q}}{\vec \eta}(t') \xrightarrow{t\to\infty}\int^t_0 dt' e^{ (t-t'){\bf Q}}{\vec \eta}(t').\label{Deltaxt}
\end{equation}
On the other hand, the deviation or fluctuation of $x_i$ about its mean value is denoted by $\delta x_i\equiv x_i- \langle x_i \rangle$, and hence  $\langle \delta x_i \rangle\equiv 0$. Since by definition $\langle \vec x\rangle =\vec X + \langle \Delta \vec x\rangle$, and from (\ref{Deltaxt}) $ \langle \Delta \vec x\rangle=0 $ to linear order, thus $\langle \vec x\rangle =\vec X +{\cal O}(\delta x ^2) $, i.e. $ \Delta{\vec x}$ and $\delta{\vec x}$ in general differ by ${\cal O}(\delta x ^2)  $. In two special cases, the local probability maximum coincides with the noisy-free fixed point, $\langle \vec x\rangle =\vec X$.
First is the case of linear  ${\vec{ F}}({\vec x}) $ (both $f_i(x_i)$ and the coupling function $h$ are linear), the resulting Langevin system is linear, then it can be easily shown that the steady-state average of ${\vec x}$ coincides with the stable fixed point: $\langle \vec x\rangle={\vec X}$ and hence  $ \delta{\vec x}=\Delta{\vec x}$. In this case, $\langle {\vec { F}}({\vec x})\rangle = {\vec { F}}(\langle{\vec x}\rangle)= {\vec { F}}({\vec X})=0$.
The second is the equilibrium case such as ${\vec{ F}}({\vec x})=-\nabla U({\vec x})$ and $\boldsymbol{\sigma}=2c{\bf I}$, the equilibrium probability distribution is Boltzmann $\psi_{eqm}\sim e^{-U({\vec x})/c}$, which implies the local peak of $\psi_{eqm} $ is ${\vec X}$.
Also for the more general equilibrium condition for noisy network dynamics, ${\bf Q}\boldsymbol{\sigma}$ is symmetric (which implies $\boldsymbol{\sigma}^{-1}{\bf Q}$ is also symmetric), then one can construct $ U({\vec x})=-\frac{\Delta{\vec x}^\intercal\boldsymbol{\sigma}^{-1}{\bf Q}\Delta{\vec x}}{2}$, which has a minimum at ${\vec X}$.

On the other hand, the interplay of nonlinearity in  ${\vec{ F}}({\vec x})$ and the non-equilibrium condition gives rise to a shift of the local maximum of $\psi_{ss}({\vec x})$ from ${\vec X}$, and one has $\langle{\vec x}\rangle={\vec X} + \langle \Delta \vec x\rangle$ with $\langle \Delta \vec x\rangle\sim {\cal O}(\delta x^2)$ as discussed above. Although $ \delta{\vec x}$ and $\Delta{\vec x}$  in general differ by $ {\cal O}(\delta x^2)$, their corresponding time-lag correlation functions (in the long-time limit for the steady-state) agree much better, as shown below:
\begin{eqnarray}
    {\bf K}_\tau&\equiv& \langle \delta{\vec x}(\tau)\delta{\vec x}^\intercal(0)\rangle \\
 &=&  \langle{\vec x}(\tau){\vec x}(0)^\intercal\rangle -  \langle{\vec x}(\tau)\rangle  \langle{\vec x}^\intercal(0)\rangle \label{ktau2} \\
  &=&  {\bf C}_\tau -  \langle\Delta{\vec x}\rangle  \langle\Delta{\vec x}^\intercal\rangle\label{ktau3}
\end{eqnarray}
where $ {\bf C}_\tau\equiv \langle\Delta{\vec x}(\tau) \Delta{\vec x}^\intercal(0)\rangle$. $ {\bf K}_\tau $ and $ {\bf C}_\tau $  differ by $ {\cal O}(\delta x^4)$, and hence  $ {\bf C}_\tau $ can be well approximated by $ {\bf K}_\tau $.

From the network reconstruction perspective, $ {\bf K}_\tau $ is measured experimentally or in simulations using (\ref{ktau2}) by recording the time-series data of $x_i(t)$. On the other hand, theoretical results relating the correlation function $ {\bf C}_\tau$ to ${\bf Q}$ and $\boldsymbol{\sigma}$ can be derived using  the linearized approximation (\ref{Deltaxt}), 
\begin{eqnarray}
  {\bf C}_\tau&=&e^{\tau{\bf Q}} {\bf C}_0\\
  \boldsymbol{\sigma}&=&-{\bf Q}{\bf C}_0-{\bf C}_0 {\bf Q}^\intercal.\label{QCCQ}
\end{eqnarray}
Adopting the approximation of  $ {\bf C}_\tau \simeq   {\bf K}_\tau $,
 one obtains the reconstruction formula 
\begin{eqnarray}
  {\bf K}_\tau&=&e^{\tau{\bf Q}} {\bf K}_0,\hbox{ and hence }  {\bf Q}=\frac{1}{\tau}\ln ( {\bf K}_\tau{\bf K}_0^{-1})\label{Ktau}\\
  \boldsymbol{\sigma}&=&-{\bf Q}{\bf K}_0-{\bf K}_0 {\bf Q}^\intercal,\label{FDRApp}
\end{eqnarray}
constituting the framework for network reconstruction for general directed networks.
(\ref{Ktau}) can be used to reconstruct ${\bf Q}$ and (\ref{FDRApp}) is the generalized fluctuation-dissipation relation for the reconstruction of the noise matrix, which is in the form of the Lyapunov equation for ${\bf K}_0$ and its (unique) solution is given by
\begin{equation}
{\bf K}_0=\int_0^\infty dt e^{t{\bf Q}} \boldsymbol{\sigma}e^{t{\bf Q}^\intercal}.\label{Lynsoln}
\end{equation}
For the equilibrium or linear ${ F}({\vec x})$ cases,  since there is no difference between  $ {\bf K}_\tau $ and $ {\bf C}_\tau $, one expects a more accurate network reconstruction performance.
In addition, ${\bf Q} $ and $ \boldsymbol{\sigma} $ can also be efficiently reconstructed from the time-series data using Stochastic force Inference.


For the equilibrium case of symmetric ${\bf Q}\boldsymbol{\sigma}$, one can show easily
\begin{equation}
  \boldsymbol{\sigma}=-2   {\bf Q} {\bf K}_0.\label{FDRsym}
  \end{equation}
Notice that the RHSs of (\ref{FDRApp}) and (\ref{FDRsym}) are the same for diagonal elements, and hence their differences will show up
only on the off-diagonal predictions of $\sigma_{ij}$. To distinguish them, one can plot the off-diagonal elements of ${\bf K}_0 {\bf Q}^\intercal $ vs.
$ {\bf Q} {\bf K}_0 $.

The fluctuations of the noisy network dynamics are time-reversal symmetric iff, i.e. $  {\bf K}_\tau= {\bf K}_\tau^\intercal $ symmetric for all $\tau$.
One can easily show that the following are equivalent conditions for time-reversal symmetric dynamics (i.e. equilibrium)
\begin{eqnarray}
     {\bf K}_\tau&=& {\bf K}_\tau^\intercal\\
     {\bf K}_0 {\bf Q}^\intercal&=& {\bf Q} {\bf K}_0=-\frac{\boldsymbol{\sigma}}{2}\\
   {\bf Q} \boldsymbol{\sigma} &=&  \boldsymbol{\sigma}{\bf Q}^\intercal.
\end{eqnarray}
\section{Detailed derivation of the conjugate force and the response function of some observables}
\subsection{General formula of the conjugate force}
\quad From the effective potential in (\ref{steady-state prob}), the conjugate force is defined as $\vec{\mathscr{F}} \equiv -\nabla_p \Phi$. Then $\nabla_p \ln{\psi_{ss}} = \delta \vec{\mathscr{F}}$. Since the effective potential $\Phi(\vec{x}; \vec{p}, \vec{X}(\vec{p}))$ depends on the fixed point and the fixed point itself depends on $\vec{p}$, one obtains
\begin{equation}
    \mathscr{F}_\alpha = -\frac{\partial \Phi}{\partial p_\alpha} - \sum_s \frac{\partial \Phi}{\partial X_s} \frac{\partial X_s}{\partial p_\alpha} = -\frac{\partial \Phi}{\partial p_\alpha} + \sum_s \frac{\partial \Phi}{\partial X_s} \left[ \mathbf{Q}^{-1} \nabla_p \vec{F} \right]_{s\alpha}.
    \label{scriptF}
\end{equation} 
Then the dependence of $\Phi$ on the fixed point is
\begin{eqnarray}
    \frac{\partial \Phi}{\partial X_s} &=& \frac{\partial }{\partial X_s} \left( \frac{1}{2} \sum_{\ell,m} ({\bf K}_0^{-1})_{\ell m} \delta x_\ell \delta x_m \right)\\
    &=&-\sum_{m}({\bf K}_0^{-1})_{sm}\delta x_m + \frac{1}{2}\sum_{\ell,m} \sum_{\mu,\nu} \frac{\partial ({\bf K}_0^{-1})_{\ell m}}{\partial Q_{\mu\nu}}\frac{\partial Q_{\mu\nu}}{\partial X_s} \delta x_\ell \delta x_m\nonumber\\
    &=& -\sum_{m}({\bf K}_0^{-1})_{sm}\delta x_m +\frac{1}{2}\sum_{\ell,m} \Bigg\{\frac{\partial ({\bf K}_0^{-1})_{\ell m}}{\partial Q_{ss}} \left[f''_s(X_s;r_s)+\sum_{n\neq s} W_{s n}\partial^2_{1}h(X_s,X_n)\right]\nonumber\\ &+&\sum_{\mu}\frac{\partial ({\bf K}_0^{-1})_{\ell m}}{\partial Q_{\mu\mu}} W_{\mu s} \partial_{1}\partial_2 h(X_\mu,X_s) + \sum_{\nu}\frac{\partial ({\bf K}_0^{-1})_{\ell m}}{\partial Q_{s\nu}} W_{s \nu}  \partial_{1}\partial_2 h(X_s,X_\nu)\nonumber\\
    &+&\sum_{\mu} \frac{\partial ({\bf K}_0^{-1})_{\ell m}}{\partial Q_{\mu s} }W_{\mu s} \partial^2_{2}h(X_\mu,X_s) \Bigg\} \delta x_\ell \delta x_m\nonumber,
    \label{pphipX}
\end{eqnarray}
where we use
\begin{eqnarray}
    \frac{\partial Q_{\mu\nu}}{\partial X_s} &=& \delta_{\mu s} \delta_{\mu\nu} \left[ f''_s(X_s;r_s)+\sum_{m\neq s} W_{sm} \partial_1^2 h(X_s,X_m) \right] + \delta_{\mu\nu} W_{\mu s} \partial_1\partial_2 h(X_\mu,X_s)\nonumber\\
    &+&\delta_{\mu s} W_{s\nu} \partial_1 \partial_2 h (X_s,X_\nu)+\delta_{\nu s} W_{\mu s} \partial_2^2 h (X_\mu,X_s).
\end{eqnarray}
$ \frac{\partial {\bf K}_0^{-1}}{\partial Q_{\mu\nu}}  $ denotes the matrix whose $i,j$ elements are $\frac{\partial ({\bf K}_0^{-1})_{ij}}{\partial Q_{\mu\nu}}$.\par
In the linearized noisy network dynamics, the steady-state probability distribution is approximated by a Gaussian distribution near its mean $\langle \vec{x} \rangle$, with $\Phi = \frac{1}{2} \delta \vec{x}^\intercal \mathbf{K}^{-1}_0 \delta \vec{x} = \frac{1}{2} (\vec{x}-\langle \vec{x} \rangle)^\intercal \mathbf{K}^{-1}_0 (\vec{x}-\langle \vec{x} \rangle)$. Note that $\langle {\vec x} \rangle$ depends on the parameters ${\vec p}$ through $\psi_{ss}$, and this dependence is generally difficult to compute explicitly in high-dimensional nonlinear Langevin systems. We therefore approximate $\langle \vec{x} \rangle$ by $\langle \vec{X} \rangle$. Nevertheless, the conjugate force ${\vec{\mathscr F}}$ and correlations such as $\langle\delta \mathbf{B} \delta\vec{\mathscr F}\rangle_{ss,{\vec p}}$ can still be computed explicitly from (\ref{scriptF}) using the functional dependence of ${\bf K}_0^{-1}$ on ${\vec p}$. From the Lyapunov equation (\ref{Lyapunov_eq}), one sees that ${\bf K}_0^{-1}={\bf K}_0^{-1}({\bf Q}({\vec p},{\vec X}({\vec p})))$. Hence the conjugate force corresponding to the parameter $p_\alpha$ can be written as
\begin{eqnarray}
{\mathscr F}_\alpha&=& - \frac{1}{2} \sum_{\ell, m}\sum_{\mu,\nu} \frac{\partial ({\bf K}_0^{-1})_{\ell m}} {\partial Q_{\mu\nu}}\frac {\partial Q_{\mu\nu}}{\partial p_\alpha} \delta x_\ell \delta x_m -\sum_{\ell, m} \left({\bf Q}^{-1} \nabla_p{\vec { F}} \Big|_{\vec X}\right)_{\ell\alpha}({\bf K}_0^{-1})_{\ell m}\delta x_m \label{Falp}\\
&+&\frac{1}{2} \sum_s \left({\bf Q}^{-1} \nabla_p{\vec { F}} \Big|_{\vec X}\right)_{s\alpha} \sum_{\ell,m} \Bigg\{ \frac{\partial (\mathbf{K}_0^{-1})_{\ell m}}{\partial Q_{ss}}\left( f_s''(X_s;r_s)+\sum_{n\neq s} W_{sn} \partial^2_1 h(X_s, X_n) \right)\nonumber\\
&+& \sum_\mu \frac{\partial (\mathbf{K}_0^{-1})_{\ell m}}{\partial Q_{\mu\mu}} W_{\mu s} \partial_1 \partial_2 h(X_\mu, X_s) + \sum_\nu \frac{\partial (\mathbf{K}_0^{-1})_{\ell m}}{\partial Q_{s\nu}} W_{s\nu}\partial_1\partial_2 h(X_s,X_\nu)\nonumber\\ &+& \sum_\mu \frac{\partial (\mathbf{K}_0^{-1})_{\ell m}}{\partial Q_{\mu s}} W_{\mu s} \partial^2_2 h(X_\mu, X_s) \Bigg\} \delta x_\ell \delta x_m .\nonumber
\end{eqnarray}
One can derive, or directly verify from (\ref{Falp}), that
\begin{equation}
    \frac{\partial X_i}{\partial p_\alpha} =\langle \delta x_i {\mathscr F}_\alpha\rangle,
\end{equation}
where $\langle \delta x_i \delta x_\ell \delta x_m \rangle=0$. By using (\ref{inverse_partialKo}), we can write down the ensemble average of the conjugate force
\begin{eqnarray}
\langle{{\mathscr F}_\alpha}\rangle&=& \frac{1}{2} \sum_{\mu,\nu} \frac {\partial Q_{\mu\nu}}{\partial p_\alpha} \text{Tr}\left({\bf K}_0^{-1} \frac{\partial {\bf K}_0} {\partial Q_{\mu\nu}}\right)-\frac{1}{2} \sum_s \left({\bf Q}^{-1} \nabla_p{\vec { F}} \Big|_{\vec X}\right)_{s\alpha}\Bigg\{ \Bigg[ f_s''(X_s;r_s)\\
&+&\sum_{n\neq s} W_{sn} \partial^2_1 h(X_s, X_n) \Bigg] \text{Tr}\left({\bf K}_0^{-1} \frac{\partial {\bf K}_0} {\partial Q_{ss}}\right)+\sum_\mu W_{\mu s} \partial_1 \partial_2 h(X_\mu, X_s) \text{Tr}\left({\bf K}_0^{-1} \frac{\partial {\bf K}_0} {\partial Q_{\mu\mu}}\right)\nonumber\\
&+& \sum_\nu W_{s\nu}\partial_1\partial_2 h(X_s,X_\nu) \text{Tr}\left({\bf K}_0^{-1} \frac{\partial {\bf K}_0} {\partial Q_{s\nu}}\right)+ \sum_\mu W_{\mu s} \partial^2_2 h(X_\mu, X_s) \text{Tr}\left({\bf K}_0^{-1} \frac{\partial {\bf K}_0} {\partial Q_{\mu s}}\right) \Bigg\}.\nonumber
\end{eqnarray}
\subsection{Conjugate force with respect to the change in the intrinsic node parameter r}
\quad Consider perturbations of an intrinsic node parameter $r_\alpha$.
\begin{equation}
    \left({\bf Q}^{-1} \nabla_r{\vec {F}}|_{\vec X} \right)_{s\alpha} =
    Q_{s\alpha}^{-1}\frac{\partial f_\alpha}{\partial r_\alpha} \bigg|_{\vec X},\quad
    \frac{\partial Q_{\mu\nu}}{\partial r_\alpha}=\frac{\partial f'_\alpha}{\partial r_\alpha}\bigg|_{\vec X}\delta_{\mu\alpha}\delta_{\mu\nu},\label{QgradF}
\end{equation}
and ${\mathscr F}_\alpha$  can then be calculated from (\ref{Falp}) to give
\begin{eqnarray}
    {\mathscr F}_\alpha&=& -\frac{1}{2} \frac{\partial f'_\alpha}{\partial r_\alpha}\Bigg|_{\vec{X}} \sum_{\ell,m} \frac{\partial (\mathbf{K}_0^{-1})_{\ell m}}{\partial Q_{\alpha\alpha}} \delta x_\ell \delta x_m - \frac{\partial f_\alpha}{\partial r_\alpha}\Bigg|_{\vec{X}} \sum_{\ell,m} Q^{-1}_{\ell\alpha} (\mathbf{K}_0^{-1})_{\ell m} \delta x_m\\
    &+&\frac{1}{2}\frac{\partial f_\alpha}{\partial r_\alpha}\Bigg|_{\vec{X}} \sum_s {Q}^{-1}_{s\alpha} \sum_{\ell,m} \Bigg\{ \frac{\partial (\mathbf{K}_0^{-1})_{\ell m}}{\partial Q_{ss}}\left( f_s''(X_s;r_s)+\sum_{n\neq s} W_{sn} \partial^2_1 h(X_s, X_n) \right)\nonumber\\
    &+& \sum_\mu \frac{\partial (\mathbf{K}_0^{-1})_{\ell m}}{\partial Q_{\mu\mu}} W_{\mu s} \partial_1 \partial_2 h(X_\mu, X_s) + \sum_\nu \frac{\partial (\mathbf{K}_0^{-1})_{\ell m}}{\partial Q_{s\nu}} W_{s\nu}\partial_1\partial_2 h(X_s,X_\nu)\nonumber\\ &+& \sum_\mu \frac{\partial (\mathbf{K}_0^{-1})_{\ell m}}{\partial Q_{\mu s}} W_{\mu s} \partial^2_2 h(X_\mu, X_s) \Bigg\} \delta x_\ell \delta x_m .\nonumber
\end{eqnarray}
Average quantities involving ${\mathscr F}_\alpha$ can then be calculated. For example, one obtains
\begin{eqnarray}
\langle{{\mathscr F}_\alpha}\rangle&=& \frac{1}{2} \frac{\partial f'_\alpha}{\partial r_\alpha}\bigg|_{\vec X} \text{Tr}\left({\bf K}_0^{-1} \frac{\partial {\bf K}_0} {\partial Q_{\alpha\alpha}}\right)-\frac{1}{2}\frac{\partial f_\alpha}{\partial r_\alpha} \bigg|_{\vec X} \sum_s Q_{s\alpha}^{-1}\Bigg\{ \Bigg[ f_s''(X_s;r_s)\\
&+&\sum_{n\neq s} W_{sn} \partial^2_1 h(X_s, X_n) \Bigg] \text{Tr}\left({\bf K}_0^{-1} \frac{\partial {\bf K}_0} {\partial Q_{ss}}\right)+\sum_\mu W_{\mu s} \partial_1 \partial_2 h(X_\mu, X_s) \text{Tr}\left({\bf K}_0^{-1} \frac{\partial {\bf K}_0} {\partial Q_{\mu\mu}}\right)\nonumber\\
&+& \sum_\nu W_{s\nu}\partial_1\partial_2 h(X_s,X_\nu) \text{Tr}\left({\bf K}_0^{-1} \frac{\partial {\bf K}_0} {\partial Q_{s\nu}}\right)+ \sum_\mu W_{\mu s} \partial^2_2 h(X_\mu, X_s) \text{Tr}\left({\bf K}_0^{-1} \frac{\partial {\bf K}_0} {\partial Q_{\mu s}}\right) \Bigg\},\nonumber
\end{eqnarray}
To calculate the response of observables that depend on $\delta \vec{x}$, one must consider multipoint correlations between $\delta x$ and $\mathscr{F}_{\alpha}$. Higher-order correlations of $\delta x_i$ can be evaluated using the cumulant expansion for a high-dimensional Gaussian distribution. In this way, quantities such as $\langle \delta x_i \delta x_j \cdots \mathscr{F}_{\alpha}\rangle$ can be computed from the corresponding moments of $\delta x$. Here we display only the two basic results, $\langle \delta x_i\mathscr{F}_\alpha\rangle$ and $\langle \delta x_i \delta x_j{\mathscr F}_\alpha\rangle$, namely
\begin{equation}
    \langle \delta x_i \mathscr{F}_{\alpha}\rangle = \frac{\partial X_i}{\partial r_\alpha},
    \label{dxdr}
\end{equation}
\begin{eqnarray}
\langle \delta x_i \delta x_j{\mathscr F}_\alpha\rangle&=& \frac{1}{2}\frac{\partial f'_\alpha}{\partial r_\alpha}\bigg|_{\vec X}\left[ ({\bf K}_0)_{ij}\text{Tr}\left({\bf K}^{-1}_0\frac{\partial {\bf K}_0}{\partial Q_{\alpha \alpha}} \right)+2 \left(\frac{\partial {\bf K}_0}{\partial Q_{\alpha \alpha}}\right)_{ij} \right]\label{dxdxscrF}\\
&-&\frac{1}{2} \frac{\partial f_\alpha}{\partial r_\alpha}\bigg|_{\vec X} \sum_s Q_{s\alpha}^{-1}\Biggl\{ \left[ f_s''(X_s;r_s)+  \sum_n W_{sn}\partial_1^2h(X_s,X_n)\right] \times\nonumber\\
& &\left[ ({\bf K}_0)_{ij} \text{Tr}\left({\bf K}^{-1}_0\frac{\partial {\bf K}_0}{\partial Q_{ss}}\right)+2 \left(\frac{\partial {\bf K}_0}{\partial Q_{ss}} \right)_{ij}\right] \nonumber\\
&+&\sum_\mu W_{\mu s}\partial_1\partial_2h(X_\mu,X_s) \left[ ({\bf K}_0)_{ij} \text{Tr} \left({\bf K}_0^{-1}\frac{\partial {\bf K}_0}{\partial Q_{\mu\mu}} \right)+2 \left(\frac{\partial {\bf K}_0}{\partial Q_{\mu\mu}} \right)_{ij}\right] \nonumber\\
&+& \sum_\nu W_{s\nu} \partial_1\partial_2h(X_s,X_\nu)  \left[ ({\bf K}_0)_{ij}\text{Tr} \left({\bf K}_0^{-1}\frac{\partial {\bf K}_0}{\partial Q_{s\nu}} \right)+2 \left(\frac{\partial {\bf K}_0}{\partial Q_{s\nu}} \right)_{ij}\right]    \nonumber\\
& +&\sum_\mu W_{\mu s}  \partial_2^2 h(X_\mu,X_s)\left[ ({\bf K}_0)_{ij} \text{Tr} \left({\bf K}_0^{-1} \frac{\partial {\bf K}_0}{\partial Q_{\mu s}} \right)+2 \left(\frac{\partial {\bf K}_0}{\partial Q_{\mu s}} \right)_{ij}\right]\Biggr\}.\nonumber
\end{eqnarray}
\subsection{Conjugate force with respect to the change in the connection weight W}
\quad We next consider perturbations of a network connection $W_{\alpha\beta}$, which directly modify the interaction topology. As in the previous subsection, we first compute the derivatives with respect to the directed connection $W_{\alpha\beta}$:
\begin{eqnarray}
\left({\bf Q}^{-1} \nabla_W{\vec { F}} |_{\vec X}\right)_{s\alpha\beta} &=& 
Q_{s\alpha}^{-1}h(X_\alpha, X_\beta)\label{QgradWF}\\
\frac{\partial Q_{\mu\nu}}{\partial W{_{\alpha\beta}}}&=& \partial_1 h(X_\alpha, X_\beta)\delta_{\mu\nu}\delta_{\mu\alpha}+\partial_2h(X_\alpha, X_\beta)\delta_{\mu\alpha}\delta_{\nu\beta}.\label{dQdW}
\end{eqnarray}
$\boldsymbol{{\mathscr F}} \equiv -\nabla_{\mathbf{W}} \Phi$ is the conjugate force associated with a change in the connection weight $W_{\alpha\beta}$, and it can be calculated from (\ref{Falp}) to give
\begin{eqnarray}
{\mathscr F}_{\alpha\beta}&=& -\frac{1}{2}\sum_{\ell,m} \biggl\{\frac{\partial ({\bf K}_0^{-1})_{\ell m}}{\partial Q_{\alpha \alpha}}\partial_1 h(X_\alpha, X_\beta)+ \frac{\partial ({\bf K}_0^{-1})_{\ell m}}{\partial Q_{\alpha \beta}}\partial_2h(X_\alpha, X_\beta)\biggr\} \delta x_\ell \delta x_m \nonumber\\
&-& h(X_\alpha, X_\beta)\sum_{\ell,m} Q^{-1}_{\ell\alpha}({\bf K}_0^{-1})_{\ell m}\delta x_m\\
&+&\frac{ h(X_\alpha, X_\beta)}{2} \sum_s Q_{s\alpha}^{-1}\sum_{\ell,m}  \Biggl\{ \frac{\partial ({\bf K}_0^{-1})_{\ell m}} {\partial Q_{ss}}\biggl[f_s''+\sum_n W_{sn}  \partial^2_1h(X_s,X_n)\biggr]\nonumber\\
&+& \sum_\mu W_{\mu s} \partial_1\partial_2h(X_\mu,X_s)\frac{\partial ({\bf K}_0^{-1})_{\ell m}} {\partial Q_{\mu\mu}} + \sum_\nu W_{s\nu}\partial_1\partial_2h(X_s,X_\nu)\frac{\partial ({\bf K}_0^{-1})_{\ell m}}{\partial Q_{s\nu}}\nonumber\\
&+& \sum_\mu  W_{\mu s}\partial^2_2h(X_\mu,X_s)\frac{\partial ({\bf K}_0^{-1})_{\ell m}}{\partial Q_{\mu s}} 
\Biggr\} \delta x_\ell \delta x_m ,\nonumber
\end{eqnarray}
and its average is given by
\begin{eqnarray}
   \langle\mathscr{F}_{\alpha\beta}\rangle  &=&-\frac{1}{2}\biggl\{ \text{Tr}\left( \frac{\partial {\bf K}_0^{-1}}{\partial Q_{\alpha \alpha}}{\bf K}_0\right) \partial_1 h(X_\alpha, X_\beta)+ \text{Tr}\left(\frac{\partial {\bf K}_0^{-1}}{\partial Q_{\alpha \beta}}{\bf K}_0\right)\partial_2h(X_\alpha, X_\beta)\biggr\} \nonumber\\
    & &+\frac{h(X_\alpha, X_\beta)}{2}\sum_s Q_{s\alpha}^{-1}  \Biggl\{
 \left[f_s''+\sum_n W_{sn}\partial_1^2h(X_s,X_n)\right] \text{Tr}\left(  \frac{\partial {\bf K}_0^{-1}}{\partial Q_{ss}}  {\bf K}_0\right) \label{avescrFab}\\
& &+ \sum_\mu W_{\mu s}\partial_1\partial_2h(X_\mu,X_s)\text{Tr}\left(  \frac{\partial {\bf K}_0^{-1}}{\partial Q_{\mu\mu}}  {\bf K}_0\right)+ \nonumber \\
 & &+ \sum_\nu  \text{Tr}\left(  \frac{\partial {\bf K}_0^{-1}}{\partial Q_{s\nu}}  {\bf K}_0\right) W_{s\nu}\partial_1\partial_2h(X_s,X_\nu) + \sum_\mu  \text{Tr}\left(  \frac{\partial {\bf K}_0^{-1}}{\partial Q_{\mu s}}  {\bf K}_0\right) W_{\mu s}\partial^2_2h(X_\mu,X_s)  
   \Biggr\}\nonumber .
\end{eqnarray}
Other relevant averages can be calculated as well:
\begin{equation}
\langle \delta x_i {\mathscr F}_{\alpha\beta}\rangle= \frac{\partial X_i}{\partial W_{\alpha\beta}} \label{dxscrFab},
\end{equation} 
\begin{eqnarray}
    \langle \delta x_i  \delta x_j{\mathscr F}_{\alpha\beta}\rangle &=&\frac{1}{2}\Bigg\{ \partial_1h(X_\alpha,X_\beta)\Bigg[ ({\bf K}_0)_{ij}\text{Tr}\Bigg({\bf K}_0^{-1} \frac{\partial {\bf K}_0}{\partial Q_{\alpha\alpha}} \Bigg) + 2\Bigg( \frac{\partial {\bf K}_0}{\partial Q_{\alpha\alpha}} \Bigg)_{ij} \Bigg]\\
    &+&\partial_2h(X_\alpha,X_\beta)\Bigg[ ({\bf K}_0)_{ij}\text{Tr}\Bigg({\bf K}_0^{-1} \frac{\partial {\bf K}_0}{\partial Q_{\alpha\beta}} \Bigg) + 2\Bigg( \frac{\partial {\bf K}_0}{\partial Q_{\alpha\beta}} \Bigg)_{ij} \Bigg]\Bigg\}\nonumber\\
    &-&\frac{1}{2}h(X_\alpha,X_\beta)\sum_s Q^{-1}_{s\alpha}\Bigg\{ f_s''\Bigg[ ({\bf K}_0)_{ij}\text{Tr}\Bigg({\bf K}_0^{-1}\frac{\partial {\bf K}_0}{\partial Q_{ss}}\Bigg)+2\Bigg( \frac{\partial {\bf K}_0}{\partial Q_{ss}} \Bigg)_{ij} \Bigg]\nonumber\\
    &+&\sum_\mu W_{\mu s}\partial_1\partial_2 h(X_\mu, X_s)\Bigg[ ({\bf K}_0)_{ij}\text{Tr}\Bigg({\bf K}_0^{-1}\frac{\partial {\bf K}_0}{\partial Q_{\mu\mu}}\Bigg)+2\Bigg( \frac{\partial {\bf K}_0}{\partial Q_{\mu\mu}} \Bigg)_{ij} \Bigg]\nonumber\\
    &+&\sum_n W_{sn}\partial_1^2 h(X_s, X_n)\Bigg[ ({\bf K}_0)_{ij}\text{Tr}\Bigg({\bf K}_0^{-1}\frac{\partial {\bf K}_0}{\partial Q_{ss}}\Bigg)+2\Bigg( \frac{\partial {\bf K}_0}{\partial Q_{ss}} \Bigg)_{ij} \Bigg]\nonumber\\
    &+&\sum_\nu W_{s\nu}\partial_1\partial_2 h(X_s, X_\nu)\Bigg[ ({\bf K}_0)_{ij}\text{Tr}\Bigg({\bf K}_0^{-1}\frac{\partial {\bf K}_0}{\partial Q_{s\nu}}\Bigg)+2\Bigg( \frac{\partial {\bf K}_0}{\partial Q_{s\nu}} \Bigg)_{ij} \Bigg]\nonumber\\
    &+&\sum_\mu W_{\mu s}\partial_2^2 h(X_\mu, X_s)\Bigg[ ({\bf K}_0)_{ij}\text{Tr}\Bigg({\bf K}_0^{-1}\frac{\partial {\bf K}_0}{\partial Q_{\mu s}}\Bigg)+2\Bigg( \frac{\partial {\bf K}_0}{\partial Q_{\mu s}} \Bigg)_{ij} \Bigg]\Bigg\}.\nonumber
\end{eqnarray}
For undirected networks, $W_{\alpha\beta}=W_{\beta\alpha}$, so changing the connection from node $\beta$ to node $\alpha$ also changes the connection from node $\alpha$ to node $\beta$ by the same amount. Hence the response function for an undirected network under a change of $W_{\alpha\beta}$ is simply the sum of the two directed-network response functions associated with the changes of $W_{\alpha\beta}$ and $W_{\beta\alpha}$.

\subsection{Conjugate force with respect to the change in the noise strength $\sigma$}
\quad Unlike perturbations of deterministic parameters, variations in the noise intensity do not shift the fixed point, but instead directly modify the steady-state fluctuations. Because the deterministic force does not depend on $\sigma$, we have
\begin{equation}
    ({\bf Q}^{-1}\nabla_\sigma \vec{F}\big|_{\vec{X}})_{s\alpha} = 0.
\end{equation}
The conjugate force associated with a change in the noise strength is then
\begin{eqnarray}
    \mathscr{F}_{\alpha\alpha} &=& -\frac{\partial \Phi}{\partial \sigma_{\alpha\alpha}} = -\frac{1}{2} \sum_{\ell, m} \frac{\partial ({\bf K}_0^{-1})_{\ell m}}{\partial \sigma_{\alpha\alpha}} \delta x_\ell \delta x_m,\\
    \langle \mathscr{F}_{\alpha\alpha} \rangle &=& \frac{1}{2} \text{Tr} \left( {\bf K}_0^{-1} \frac{\partial {\bf K}_0}{\partial \sigma_{\alpha\alpha}} \right).
\end{eqnarray}
The fixed point does not change when the noise strengths are perturbed:
\begin{equation}
    \langle \delta x_i \mathscr{F}_{\alpha\alpha}\rangle = \frac{\partial X_i}{\partial \sigma_{\alpha\alpha}}=0.
\end{equation}
The two-point correlation is
\begin{equation}
    \langle \delta x_i \delta x_j \mathscr{F}_{\alpha\alpha} \rangle = \frac{1}{2} \left[ (\mathbf{K}_0)_{ij} \text{Tr}\left( \mathbf{K}_0^{-1}\frac{\partial \mathbf{K}_0}{\partial \sigma_{\alpha\alpha}} \right) + 2 \left( \frac{\partial \mathbf{K}_0}{\partial \sigma_{\alpha\alpha}} \right)_{ij} \right]
\end{equation}

\subsection{General formula of the time-dependent conjugate force}
\quad For time-dependent perturbations, one needs the time-lagged correlation function between the observable and the conjugate force. To calculate the time-dependent change of the fixed point, in analogy with (\ref{scriptF}), we consider $\langle \delta \vec{x}(\tau) \vec{\mathscr{F}}^\intercal(0) \rangle$. Here $\tau$ denotes the time lag appearing in the correlation functions, while $t$ is reserved for the time variable in the response functions. For the linearized dynamics,
\begin{equation}
    \delta \vec{x}(\tau)=e^{\tau \mathbf{Q}} \delta \vec{x}(0)+ \int_0^\tau ds\, e^{(\tau-s)\mathbf{Q}}\vec{\eta}(s).
\end{equation}
Therefore
\begin{eqnarray}
    \langle \delta x_i(\tau) {\mathscr F}_\alpha(0)\rangle &=& \sum_{\ell} \left( e^{\tau{\bf Q}} \right)_{i\ell}\langle \delta x_\ell(0) {\mathscr F}_\alpha(0)\rangle \nonumber\\
    &+&\int_0^\tau ds\, \sum_k \left( e^{(\tau-s){\bf Q}} \right)_{ik}\langle \eta_k(s) {\mathscr F}_\alpha(0)\rangle.
\end{eqnarray}
Since $s>0$ in the noise integral and ${\mathscr F}_\alpha(0)$ depends only on the state at time $0$, the future noise is uncorrelated with ${\mathscr F}_\alpha(0)$, so the second term vanishes. Using the equal-time identity $\langle \delta x_\ell {\mathscr F}_\alpha \rangle=\partial X_\ell/\partial p_\alpha$, one obtains
\begin{equation}
    \langle \delta x_i(\tau) {\mathscr F}_\alpha(0)\rangle= \sum_{\ell} \frac{\partial X_\ell}{\partial p_\alpha}
\left( e^{\tau{\bf Q}} \right)_{i\ell}, \label{dxtauscrF}
\end{equation}
\begin{eqnarray}
\langle \delta x_i(\tau) \delta x_j(\tau){\mathscr F}_\alpha(0)\rangle&=&  \Biggl\{- \frac{1}{2} \sum_{m,\ell} \sum_{\mu,\nu} \frac{\partial ({\bf K}_0^{-1})_{m\ell}} {\partial Q_{\mu\nu}}\frac {\partial Q_{\mu\nu}}{\partial p_\alpha}\\
& &+\frac{1}{2}  \sum_s \left({\bf Q}^{-1} \nabla_p{\vec { F}} |_{\vec X}\right)_{s\alpha} \sum_{m,\ell} \Bigg[ [f_s''+\sum_n W_{sn}  \partial^2_1 h(X_s,X_n)] \frac{\partial ({\bf K}_0^{-1})_{m\ell}} {\partial Q_{ss}}\nonumber\\
 & & + \sum_\mu W_{\mu s} \partial_1\partial_2h(X_\mu,X_s)\frac{\partial ({\bf K}_0^{-1})_{m\ell}} {\partial Q_{\mu\mu}}+ \sum_\nu W_{s\nu}\partial_1\partial_2h(X_s,X_\nu)\frac{\partial ({\bf K}_0^{-1})_{m\ell}}{\partial Q_{s\nu}}  \nonumber \\
 & &+ \sum_\mu W_{\mu s}\partial^2_2h(X_\mu,X_s) \frac{\partial ({\bf K}_0^{-1})_{m\ell}}{\partial Q_{\mu s}} \Bigg]
\Biggr\} \langle \delta x_i(\tau) \delta x_j(\tau)    \delta x_\ell(0) \delta x_m(0) \rangle \nonumber\\
  &=&- \frac{1}{2}  \sum_{\mu,\nu}\frac {\partial Q_{\mu\nu}}{\partial p_\alpha}\biggl[ ({\bf K}_0)_{ij}\text{Tr}\left(\frac{\partial {\bf K}_0^{-1}}{\partial Q_{\mu \nu}} {\bf K}_0\right)+2 \left({\bf K}_\tau\frac{\partial {\bf K}_0^{-1}}{\partial Q_{\mu \nu}} {\bf K}_\tau^\intercal\right)_{ij}\bigg]\nonumber\\
  &&+\frac{1}{2}  \sum_s \left({\bf Q}^{-1} \nabla_p{\vec { F}} |_{\vec X}\right)_{s\alpha}  \Biggl\{ [ f_s''+  \sum_n W_{sn}\partial_1^2h(X_s,X_n)] \times\nonumber\\
    & &\biggl[ ({\bf K}_0)_{ij}\text{Tr}\left(\frac{\partial {\bf K}_0^{-1}}{\partial Q_{ss}} {\bf K}_0\right)+2 \left({\bf K}_\tau\frac{\partial {\bf K}_0^{-1}}{\partial Q_{ss}} {\bf K}_\tau^\intercal\right)_{ij}\bigg] \label{dxdxtauscrF}\\
   & & +\sum_\mu W_{\mu s}\partial_1\partial_2h(X_\mu,X_s) \biggl[ ({\bf K}_0)_{ij}\text{Tr}\left(\frac{\partial {\bf K}_0^{-1}}{\partial Q_{\mu\mu}} {\bf K}_0\right)+2 \left({\bf K}_\tau\frac{\partial {\bf K}_0^{-1}}{\partial Q_{\mu\mu}} {\bf K}_\tau^\intercal\right)_{ij}\bigg] \nonumber\\
   & & +  \sum_\nu W_{s\nu} \partial_1\partial_2h(X_s,X_\nu)  \biggl[ ({\bf K}_0)_{ij}\text{Tr}\left(\frac{\partial {\bf K}_0^{-1}}{\partial Q_{s\nu}} {\bf K}_0\right)+2 \left({\bf K}_\tau\frac{\partial {\bf K}_0^{-1}}{\partial Q_{s\nu}} {\bf K}_\tau^\intercal\right)_{ij}\bigg]    \nonumber\\
   & & +\sum_\mu W_{\mu s}  \partial_2^2h(X_\mu,X_s)\biggl[ ({\bf K}_0)_{ij}\text{Tr}\left(\frac{\partial {\bf K}_0^{-1}}{\partial Q_{\mu s}} {\bf K}_0\right)+2 \left({\bf K}_\tau\frac{\partial {\bf K}_0^{-1}}{\partial Q_{\mu s}} {\bf K}_\tau^\intercal\right)_{ij}\bigg]
   \Biggr\} \nonumber
\end{eqnarray}
where 
\begin{eqnarray}
\langle \delta x_i(\tau) \delta x_j (\tau)  \delta x_\ell(0) \delta x_m (0)\rangle
&=&({\bf K}_0)_{ij}({\bf K}_0)_{\ell m} + ({\bf K}_\tau)_{im}({\bf K}_\tau)_{j\ell }+ ({\bf K}_\tau)_{i\ell}({\bf K}_\tau)_{jm} \label{dx4tau}
\end{eqnarray}
which is used in the last equality. Note that ${\bf K}_\tau= e^{\tau{\bf Q}}{\bf K}_0$, and the right-hand side of (\ref{dxdxtauscrF}) can therefore be calculated explicitly in terms of the network model parameters.
The time-dependent response function of the state is
\begin{eqnarray}
    \chi_{i\alpha}(t) &=& \chi_{i\alpha}^{(ss)}-\langle \delta x_i(t) \mathscr{F}_\alpha(0) \rangle\\
    &=& \frac{\partial X_i}{\partial p_\alpha}
-\sum_{\ell} \frac{\partial X_\ell}{\partial p_\alpha}
\left( e^{t{\bf Q}} \right)_{i\ell}\nonumber
\end{eqnarray}
The non-steady response function of the force is
\begin{equation}
    \chi_{i\alpha}(t)= \chi_{i\alpha}^{(ss)}-\langle \delta F_i(t) {\mathscr F}_\alpha(0)\rangle
\end{equation}
\begin{eqnarray}
\langle \delta {F}_i(t) \mathscr F_\alpha(0)\rangle&=& \langle \delta f_i(t)  {\mathscr F}_\alpha(0)\rangle +\sum_{j\neq i} W_{ij} \Biggl\{ \partial_1 h(X_i,X_j) \langle \delta x_i(t) {\mathscr F}_\alpha(0)\rangle \nonumber\\ &+&\partial_2h(X_i,X_j)\langle \delta x_j(t){\mathscr F}_\alpha(0)\rangle 
+\frac{1}{2} \Big[\partial_1^2h(X_i,X_j) \langle \delta x_i^2(t) {\mathscr F}_\alpha(0) \rangle\nonumber\\
&+& \partial_2^2h(X_i,X_j) \langle \delta x_j^2(t) {\mathscr F}_\alpha(0)\rangle \Big] +\partial_1 \partial_2h(X_i,X_j)\langle \delta x_i(t) \delta x_j(t) {\mathscr F}_\alpha(0)\rangle \Biggr\}.
\end{eqnarray}
If we only consider the linear term, then it becomes
\begin{equation}
    \langle \delta F_i(t) \mathscr{F}_\alpha(0) \rangle = \sum_{\ell} ({\bf Q}e^{t{\bf Q}})_{i\ell} \frac{\partial X_\ell}{\partial p_\alpha}.
\end{equation}
The non-steady, time-dependent response function of the fluctuation correlation can then be written as
\begin{equation}
    \chi_{ij\alpha}(t)= \chi_{ij\alpha}^{(ss)}-[\langle \delta x_i(t) \delta x_j(t){\mathscr F}_\alpha(0)\rangle -({\bf K}_0)_{ij}\langle{\mathscr F}_\alpha\rangle ],\hbox{  where }
\end{equation}
\begin{eqnarray}
\langle \delta x_i(t) \delta x_j(t){\mathscr F}_\alpha(0)\rangle -({\bf K}_0)_{ij}\langle{\mathscr F}_\alpha\rangle&=&  \sum_{\mu,\nu}\frac {\partial Q_{\mu\nu}}{\partial p_\alpha}\left(e^{t{\bf Q}}\frac{\partial {\bf K}_0}{\partial Q_{\mu \nu}}e^{t{\bf Q}^\intercal}\right)_{ij}-  \sum_s \left({\bf Q}^{-1} \nabla_p{\vec { F}}|_{\vec X} \right)_{s\alpha} \times\nonumber\\
& &\Biggl\{ [ f_s''+  \sum_n W_{sn}\partial_1^2h(X_s,X_n)] \left(e^{t{\bf Q}}\frac{\partial {\bf K}_0}{\partial Q_{ss}}e^{t{\bf Q}^\intercal}\right)_{ij} \\
   & & +\sum_\mu W_{\mu s}\partial_1\partial_2h(X_\mu,X_s) \left(e^{t{\bf Q}}\frac{\partial {\bf K}_0}{\partial Q_{\mu\mu}}e^{t{\bf Q}^\intercal}\right)_{ij}\nonumber\\
   & & +  \sum_\nu W_{s\nu} \partial_1\partial_2h(X_s,X_\nu)    \left(e^{t{\bf Q}}\frac{\partial {\bf K}_0}{\partial Q_{s\nu}}e^{t{\bf Q}^\intercal}\right)_{ij}    \nonumber\\
   & & +\sum_\mu W_{\mu s}  \partial_2^2h(X_\mu,X_s)\left(e^{t{\bf Q}}\frac{\partial {\bf K}_0}{\partial Q_{\mu s}}e^{t{\bf Q}^\intercal}\right)_{ij}
   \Biggr\}. \nonumber
\end{eqnarray}

\section{Spectral solution of the Lyapunov equation and transient protocol integrals}
\quad This appendix collects two technical steps used in the main text: the spectral solution of the Lyapunov equation and the mode expansion of the time-dependent convolution integrals in Sec. \ref{sec:time-dependent}.
\subsection*{Spectral solution of the Lyapunov equation}
Consider the Lyapunov equation for the unknown matrix ${\bf K}$ \cite{Behr2019SylvesterLyapunov}
\begin{equation}
{\bf Q}{\bf K}+{\bf K}{\bf Q}^\intercal={\bf S},
\label{lyap}
\end{equation}
where ${\bf Q}$ and ${\bf S}$ are given matrices and ${\bf S}$ is symmetric.
Assume that ${\bf Q}$ is diagonalizable:
\begin{equation}
{\bf Q}={\bf P}\boldsymbol{\Lambda}{\bf P}^{-1},\quad
{\bf P}=({\vec u_1},{\vec u_2},\cdots,{\vec u_N}),\quad
\boldsymbol{\Lambda}=\hbox{diag}(\lambda_1,\lambda_2,\cdots,\lambda_N),
\end{equation}
with ${\bf Q}{\vec u_i}=\lambda_i{\vec u_i}$. Let the columns of
${\bf P}^{-\intercal}$ be the left eigenvectors,
\begin{equation}
{\bf P}^{-\intercal}=({\vec v_1},{\vec v_2},\cdots,{\vec v_N}),
\end{equation}
so that
\begin{equation}
{\vec v_i}^\intercal{\vec u_j}=\delta_{ij},\qquad
\sum_{j=1}^N {\vec u_j}{\vec v_j}^\intercal={\bf I}.
\end{equation}
Defining ${\tilde{\bf K}}={\bf P}^{-1}{\bf K}{\bf P}^{-\intercal}$, Eq. (\ref{lyap}) becomes
\begin{eqnarray}
\boldsymbol{\Lambda}{\tilde{\bf K}}+{\tilde{\bf K}}\boldsymbol{\Lambda}
&=&{\bf P}^{-1}{\bf S}{\bf P}^{-\intercal},\\
{\tilde{\bf K}}_{ij}
&=&\frac{({\bf P}^{-1}{\bf S}{\bf P}^{-\intercal})_{ij}}{\lambda_i+\lambda_j}.
\label{Ktil}
\end{eqnarray}
Hence
\begin{equation}
{\bf K}={\bf P}{\tilde{\bf K}}{\bf P}^\intercal.
\end{equation}
For the stable fixed points considered here, $\mathrm{Re}\,\lambda_i<0$, so $\lambda_i+\lambda_j\neq 0$ in Eq. (\ref{Ktil}). Using the biorthogonal decomposition
\begin{equation}
{\bf Q}=\sum_{m=1}^N \lambda_m {\vec u_m}{\vec v_m}^\intercal,
\end{equation}
the same solution can be written more compactly as
\begin{equation}
{\bf K}=\sum_{m,\ell=1}^N
\frac{({\vec v_m}^\intercal{\bf S}{\vec v_\ell})\,{\vec u_m}{\vec u_\ell}^\intercal}
{\lambda_m+\lambda_\ell}.
\end{equation}
\subsection*{Mode expansion for the time-dependent protocol integrals}
The time-dependent formulas in Sec. \ref{sec:time-dependent} follow from the decomposition of $\mathbf{Q}$. Since
\begin{equation}
e^{t{\bf Q}}=\sum_{k=1}^N e^{\lambda_k t}{\vec u_k}{\vec v_k}^\intercal,
\qquad
{\bf Q}^{-1}=\sum_{k=1}^N \frac{{\vec u_k}{\vec v_k}^\intercal}{\lambda_k},
\end{equation}
Eq. (\ref{xi(t)drlinear}) becomes
\begin{equation}
\langle x_i(t)\rangle_{r+\delta r_\alpha(t)}=X_i+(a_\alpha-X_\alpha)\bigg[
\sum_{k=1}^N \frac{u_{ik}v_{\alpha k}}{\lambda_k}\int_0^t dt'
e^{(t-t')\lambda_k}\delta{\dot r}_\alpha(t')
-Q^{-1}_{i\alpha}\delta r_\alpha(t)\bigg],
\label{xi(t)drlinearb}
\end{equation}
where $u_{ik}\equiv({\vec u_k})_i$ and $v_{\alpha k}\equiv({\vec v_k})_\alpha$.
The time-dependent relaxation response is then
\begin{equation}
      \langle x_i(t)\rangle=\begin{cases}  X_i+ \frac{b}{2} (a_\alpha-X_\alpha)\bigg[\sum_{k=1}^N\frac{u_{ik}v_{\alpha k}}{\lambda_k(1+(\frac{\lambda_k\tau}{\pi})^2)}\left(e^{\lambda_kt}-\cos\frac{\pi t}{\tau}-\frac{\lambda_k\tau}{\pi}\sin\frac{\pi t}{\tau}  \right)-&Q^{-1}_{i\alpha} (1-\cos\frac{\pi t}{\tau})\bigg] \\
     & 0\leq t\leq\tau\\
     X_i-b (a_\alpha-X_\alpha)Q^{-1}_{i\alpha}+\frac{b}{2} (a_\alpha-X_\alpha)\sum_{k=1}^N\frac{u_{ik}v_{\alpha k}[ e^{\lambda_kt}+e^{\lambda_k(t-\tau)}]}{\lambda_k(1+(\frac{\lambda_k\tau}{\pi})^2)} &t>\tau
     \end{cases}.
\end{equation}
For the force response,
\begin{equation}
\langle { F}_i(t)\rangle = \begin{cases}  (a_\alpha-X_\alpha)\sum_{k=1}^N u_{ik}v_{\alpha k}\int_0^tdt' \delta {\dot r}_\alpha(t')e^{(t-t')\lambda_k}, &0\leq t\leq\tau\\
(a_\alpha-X_\alpha)\sum_{k=1}^N u_{ik}v_{\alpha k}\int_0^\tau dt' \delta {\dot r}_\alpha(t')e^{(t-t')\lambda_k}, & t>\tau
\end{cases}
\end{equation}
\begin{equation}
\langle { F}_i(t)\rangle= \begin{cases} \frac{b}{2}(a_\alpha-X_\alpha)\sum_{k=1}^N \frac{u_{ik}v_{\alpha k}}{1+\left(\frac{\lambda_k\tau}{\pi}\right)^2}\left[ e^{\lambda_k t} - \cos\left( \frac{\pi t}{\tau} \right) - \frac{\lambda_k \tau}{\pi}\sin\left( \frac{\pi t}{\tau} \right) \right], &0\leq t\leq\tau\\
\frac{b}{2}(a_\alpha-X_\alpha)\sum_{k=1}^N \frac{u_{ik}v_{\alpha k}}{1+\left(\frac{\lambda_k\tau}{\pi}\right)^2} \left[ e^{\lambda_k t}+e^{\lambda_k(t-\tau)} \right], & t>\tau
\end{cases}.
\end{equation}
Likewise, for the covariance response,
\begin{eqnarray}
&&\int_0^t dt' \left(
e^{(t-t'){\bf Q}}
\frac{\partial {\bf K}_0}{\partial Q_{\alpha\alpha}}
e^{(t-t'){\bf Q}^\intercal}
\right)_{ij}\delta \dot{r}_\alpha(t') \nonumber\\
&=&\sum_{m,\ell=1}^N
u_{im}u_{j\ell}\left(
{\vec v_m}^\intercal
\frac{\partial {\bf K}_0}{\partial Q_{\alpha\alpha}}
{\vec v_\ell}
\right)I_{m\ell}(t),
\end{eqnarray}
with
\begin{equation}
I_{m\ell}(t)\equiv
\int_0^t dt' \, e^{(t-t')(\lambda_m+\lambda_\ell)}\delta \dot r_\alpha(t')
=\frac{b\pi}{2\tau}\int_0^{\min(t,\tau)} dt' \, e^{(t-t')(\lambda_m+\lambda_\ell)}\sin\left(\frac{\pi t'}{\tau}\right),
\label{Iml}
\end{equation}
which evaluates to
\begin{equation}
I_{m\ell}(t)=
\begin{cases}
\displaystyle
\frac{b}{2\left[1+\left(\frac{(\lambda_m+\lambda_\ell)\tau}{\pi}\right)^2\right]}
\left[
e^{(\lambda_m+\lambda_\ell)t}
-\cos\left(\frac{\pi t}{\tau}\right)
-\frac{(\lambda_m+\lambda_\ell)\tau}{\pi}\sin\left(\frac{\pi t}{\tau}\right)
\right], & 0\le t\le \tau, \\[1.2ex]
\displaystyle
\frac{b}{2\left[1+\left(\frac{(\lambda_m+\lambda_\ell)\tau}{\pi}\right)^2\right]}
\left[
e^{(\lambda_m+\lambda_\ell)t}
+e^{(\lambda_m+\lambda_\ell)(t-\tau)}
\right], & t>\tau .
\end{cases}
\end{equation}
The edge-perturbation formulas in Sec. \ref{sec:time-dependent} follow from the same expansion, with only the prefactor and the covariance derivative replaced according to the perturbation under consideration.

The corresponding time-dependent response of each node state variable is
\begin{equation}
      \langle x_i(t)\rangle=\begin{cases}  X_i+ \frac{b}{2} (X_\beta-X_\alpha)\bigg[\sum_{k=1}^N\frac{u_{ik}v_{\alpha k}}{\lambda_k(1+(\frac{\lambda_k\tau}{\pi})^2)}\left(e^{\lambda_kt}-\cos\frac{\pi t}{\tau}-\frac{\lambda_k\tau}{\pi}\sin\frac{\pi t}{\tau}  \right)&-Q^{-1}_{i\alpha} \left(1-\cos\frac{\pi t}{\tau}\right)\bigg] \\
     & 0\leq t\leq\tau\\
     X_i-b (X_\beta-X_\alpha)Q^{-1}_{i\alpha}+\frac{b}{2} (X_\beta-X_\alpha)\sum_{k=1}^N\frac{u_{ik}v_{\alpha k}[ e^{\lambda_kt}+e^{\lambda_k(t-\tau)}]}{\lambda_k(1+(\frac{\lambda_k\tau}{\pi})^2)} &t>\tau
     \end{cases}.
\end{equation}

\begin{equation}
\langle { F}_i(t)\rangle= \begin{cases} \frac{b}{2}(X_\beta-X_\alpha)\sum_{k=1}^N \frac{u_{ik}v_{\alpha k}}{1+\left(\frac{\lambda_k\tau}{\pi}\right)^2}\left[ e^{\lambda_k t} - \cos\left( \frac{\pi t}{\tau} \right) - \frac{\lambda_k \tau}{\pi}\sin\left( \frac{\pi t}{\tau} \right) \right], &0\leq t\leq\tau\\
\frac{b}{2}(X_\beta-X_\alpha)\sum_{k=1}^N \frac{u_{ik}v_{\alpha k}}{1+\left(\frac{\lambda_k\tau}{\pi}\right)^2} \left[ e^{\lambda_k t}+e^{\lambda_k(t-\tau)} \right], & t>\tau
\end{cases}.
\end{equation}

The time-dependent response function of the covariance is
\begin{eqnarray}
    {\bf K}_{ij}(t,t) &=& ({\bf K_0})_{ij} + \chi_{ij\alpha\beta}^{ss} \delta W_{\alpha\beta}(t)+\int_0^t dt' \left[ e^{(t-t'){\bf Q}}\left(-\frac{\partial {\bf K}_0}{\partial Q_{\alpha\alpha}} +\frac{\partial {\bf K}_0}{\partial Q_{\alpha\beta}}\right) e^{(t-t'){\bf Q}^\intercal} \right]_{ij} \delta \dot{W}_{\alpha\beta}(t')\nonumber\\
    &=&({\bf K_0})_{ij}+\left[- \frac{\partial {\bf K_0}}{\partial Q_{\alpha\alpha}}+\frac{\partial {\bf K_0}}{\partial Q_{\alpha\beta}} \right]_{ij}\delta W_{\alpha\beta}(t)\nonumber\\
    &+& \sum_{m,\ell} u_{im}u_{j\ell} \left[ \vec{v}^\intercal_m \left( \frac{\partial {\bf K_0}}{\partial Q_{\alpha\alpha}}-\frac{\partial {\bf K_0}}{\partial Q_{\alpha\beta}} \right) \vec{v}_\ell \right] I_{m\ell}(t).
\end{eqnarray}


\bibliographystyle{apsrev}
\bibliography{ref}

\end{document}
